% concatenated from main.tex
\documentclass[a4paper]{article}

\usepackage{fullpage}
\usepackage{comment}

\usepackage{xcolor}

\usepackage{amsthm,amsmath,amssymb,mathtools}
\usepackage[hidelinks]{hyperref}
\usepackage[nameinlink]{cleveref}

\usepackage{enumitem}

\newtheorem{theorem}{Theorem}[section]
\newtheorem{lemma}[theorem]{Lemma}
\newtheorem{corollary}[theorem]{Corollary}

\theoremstyle{definition}
\newtheorem{remark}[theorem]{Remark}

\usepackage{todonotes}

\newcommand{\Rd}{{\mathbb{R}^d}}
\newcommand{\dRL}{{^{RL}\partial_t^\alpha}}

\newcommand{\asympint}[2]{\overset {(#1,#2)}\asymp }

\usepackage[backend=biber,url=false,giveninits,maxbibnames=10]{biblatex}
\addbibresource{references.bib}
\AtEveryBibitem{%
  \clearfield{isbn}
  \clearfield{issn}
  \clearfield{month}
  \clearfield{series}
  \clearfield{urldate}
  \clearfield{pubstate}
  \clearfield{url}

  \ifentrytype{book}{}{%
    \clearlist{publisher}
    \clearname{editor}
  }
}

\DeclareMathOperator{\fixedpoint}{\mathfrak K}

\usepackage{authblk}

%
\renewcommand\Affilfont{\normalsize}

\title{Self-similar solutions to the \\
 time-fractional Porous-Medium Equation}
\author[1,2]{David Gómez-Castro}
\author[3]{\L{}ukasz P\l{}ociniczak}
\author[1]{Juan Luis Vázquez}
\affil[1]{Universidad Autónoma de Madrid}
\affil[2]{Instituto de Ciencias Matemáticas, CSIC}
\affil[3]{Wrocław University of Science and Technology}

%
\date{}

\begin{document}

\maketitle

\begin{abstract}
    We show the existence of self-similar solutions with constant finite mass to the time-fractional Porous-Medium Equation for all spatial dimensions $d \ge 1$ and all exponents $m>m_c=(d-2)_+/d$. This range is optimal. We find two types of solution depending on the exponent: compactly supported solutions in the slow-diffusion range $m > 1$ and positive solutions with heavy tails in the sub-critical fast-diffusion range $m_c < m < 1$. The self-similar solutions in the linear case $m=1$ were already known explicitly obtained by the Fourier transform, and we discuss their properties in our settings and the limit $m \to 1$.

    \noindent \textbf{Keywords.} Porous medium equation; fast diffusion equation;  Caputo fractional-time derivative; self-similar solutions.
    %
    %
\end{abstract}

%
%

%
%

\hfill {\sl Dedicated to the memory of G. I. Barenblatt}


%
\section{Introduction and main results}
In this paper, we discuss the nonnegative self-similar solutions of the nonlinear time-fractional diffusion equation for $m > 0$ and $\alpha \in (0,1)$ given by
\begin{equation}
  \label{eq:main}
  \partial_t^\alpha u = \Delta (u^{m})
\end{equation}
for $t > 0$ and $x \in \mathbb R^d$ for any $d \ge 1$.
Here $\partial_t^\alpha$ is the Caputo derivative
\begin{equation}
\label{eq:Caputo}
  \partial_t^\alpha u(t,x) = \int_0^t \frac{(t-s)^{-\alpha}}{\Gamma(1-\alpha)} \partial_t u(s,x) ds.
\end{equation}
The solutions start from a Dirac delta $u(0,\cdot) = M \delta_0$ in the sense of initial traces:
\begin{equation}
  \label{eq:initial Dirac}
  \operatorname*{ess\,lim}_{t \to 0^+} \int_{\Rd} u(t,x) \varphi(x) dx = M\varphi(0),
\end{equation}
for all $\varphi \in C^\infty_c (\Rd)$. We look for nonnegative self-similar solutions $u(t,x)$ of the equation that have a finite mass for all $t>0$, and, moreover, they conserve for all times the initial mass $M>0$ in the sense that
\begin{equation}
  \label{eq:mass conservation}
  \int_{\mathbb R^d} u(t,x) dx = M \qquad \text{for all } t > 0.
\end{equation}
By self-similar solution, we refer to space-time function, which can be re-scaled to a profile $U$ in the sense that
\begin{equation}\label{eq:ssf}
  u(t,x) = t^{-a} U(t^{-b} |x|),
\end{equation}
for similarity exponents $a,b$ to be appropriately chosen.
Our main contribution is the study of the existence and key properties of $U$ and its dependence on parameters $\alpha,m,d,$ and mass $M$.

\medskip

The classical case $\alpha = 1$ was the source of a significant body of work in the twentieth century. The linear diffusion case $m = 1$ is the most famous, where $u$ is the heat kernel and $u$ is the Gaussian profile. The slow-diffusion $m > 1$ solutions lead to the compactly-supported Barenblatt solutions \cite{MR1426127} (see also \cite{MR46217,zel1950towards}), which started the study of the theory of free boundaries in non-linear diffusion. 
These naturally arise in many physical contexts, such as gas flow through heterogeneous media, water percolation in the soil \cite{bear2013dynamics}, or even in biology \cite{vafai2010porous} and medicine \cite{lavrova2021barenblatt}. The nonlinear character of the diffusivity helps to model such physically important features as finite speed of propagation and sharp interfaces. The former is a central characteristic of fluid flow through porous media where disturbances do not propagate with an infinite speed (as modeled by the classical heat equation). Moreover, Barenblatt profiles frequently act as attractors to which solutions starting from general initial conditions converge asymptotically \cite{vazquez2014barenblatt}. \normalcolor In the fast-diffusion case $0 < m < 1$ we have self-similar power-type solutions. In this setting, one must distinguish the mildly-singular fast-diffusion case $m_c < m < 1$, where $m_c = (1 - 2/d)_+$, in which finite-mass self-similar solutions exist. In the very-singular fast-diffusion $0 < m < m_c$ the self-similar solutions have infinite mass. In the very-singular case, there is still a theory of backward-in-time self-similarity, which we will not discuss here. 
We also discuss the mesa problem, i.e. the limit as $m\to \infty$, which is well-understood in the classical case $\alpha = 1$ \cite{MR983523,MR883709}.
For details on the Porous-Medium equation, we refer the reader to \cite{Vazquez2006,Vazquez2006a}.
For a general discussion on self-similarity and mass conservation see \cite{vazquez2025surveymassconservationselfsimilarity}.

\medskip

The theory of existence and uniqueness and key properties for the fractional case $\alpha \in (0,1)$ is an active field of research.
The case of linear diffusion, that is, for $m = 1$, the existence was proven in \cite{Zacher2009}, and other properties in \cite{vergara2015,dier2020,cortazarHeatEquationMemory2021}. 
For $m > 1$ the closest literature that deals with the existence and properties is
\cite{djidaNonlocalTimeporousMedium2019,DjidaNietoArea2019}
for the nonlocal pressure case where $\Delta u^m$ is replaced by $\operatorname{div}(m u^{m-1} \nabla p)$ and $(-\Delta)^s p = u$. Our problem represents the $s = 0$ case. 

\medskip


The theory of self-similar solutions for the fractional-in-time problem with linear diffusion $m = 1$ is known \cite{MR1061448, MR2047909,MR3631303} where the Fourier transform of the self-similar solutions is of closed form given by Mittag-Leffler functions (see \eqref{eq:m=1 self-similar}). 
The self-similar solutions for $m > 1$ and $d=1$ were studied in \cite{CaballeroOkrasinska-PlociniczakPlociniczakSadarangani2025} by arriving at an integral equation (see \eqref{eq:profile equation} below) and applying Schauder fixed-point arguments. 
A numerical scheme for $m > 1$ and $d=1$ is presented in \cite{Plociniczak2019}.

\medskip

There are several other different directions of problems with fractional-in-time diffusion to cite a few:
with fractional pressure in $\mathbb R^d$ we have \cite{allen2016}.
The study of the case $m > 1$ in bounded domains is done in \cite{bonforteTimeFractionalPorousMedium2024}.
In addition, the abstract approach of flows in the Hilbert space was undertaken in \cite{akagi2019fractional} where both the linear and nonlinear versions of the main equation have been considered (along with other generalizations). 
Also, see the latest developments in \cite{akagi2025time}. 
Moreover, several existence and uniqueness results for a general nonlinear parabolic problem with the Caputo derivative have been proved in \cite{topp2017existence} via the viscosity solution framework. 

\paragraph{Outline of main results.} 
We provide a complete study of self-similar profiles for $d \ge 1$, $\alpha \in (0,1)$, $m > m_c = (\tfrac{d-2}d)_+$ that extends the previous results for $m = 1$ in \cite{MR1829592,dier2020}, and $m > 1$ and $d=1$ in \cite{CaballeroOkrasinska-PlociniczakPlociniczakSadarangani2025}.
Similarly to \cite{CaballeroOkrasinska-PlociniczakPlociniczakSadarangani2025} we arrive at an integral equation \eqref{eq:profile equation}, where we have significantly simplified the expression of the kernel in \eqref{eq:Q} even in dimension $d=1$, and we solve it using the method of sub and super-solutions in all cases.
%
We can summarize the results presented later in this section as follows:
\begin{itemize}
  \item \textbf{Existence} of profiles for all $m>m_c$.

    \item \textbf{Uniqueness} of profiles for $m\ge 1$.
    
  \item \textbf{Notion of weak solution} satisfied by the self-similar solution.
  
  \item \textbf{Sharp behaviour at $z=0$.}
  For all $m > m_c$ we see that $U$ is bounded in $d=1$ and $U^m$ behaves like the Newtonian potential for $d \ge 2$. The critical exponent $m_c$ is optimal. 
  
    \item \textbf{Compact support for $m > 1$} where solutions behave like a simple power near the free bundary. The sharp constant is provided.

    \item \textbf{Tail behavior for $m \in (m_c,1)$} 
    The solutions have power-like tails when $m_c < m < 1$, which can easily be computed by scaling, and they form an asymptotic fan around a very-singular solution of infinite mass.
    The exponential tail of the linear case $m=1$ is quite complicated. On the contrary, for $m< 1$ the tail decays like a simple power.
    We match these cases by proving the limit $m \to 1$.

  \item \textbf{Mesa limit $m\rightarrow \infty$} where the solutions are multiples of the characteristic function $\chi_{B_1}$. Curiously, this mesa-like behaviour coincides exactly with the one found for $\alpha=1$.

    \item \textbf{Classical limit $\alpha \to 1$} where we recover the usual Barenblatt profiles.
  
  \item \textbf{The numerical scheme} for the case $m>1$ is constructed by Picard's iterations. We illustrate the behaviors with respect to $\alpha$ and $m$ in \Cref{fig:limit alpha->1,fig:mesa}. 
\end{itemize}
\normalcolor

\paragraph{Notation.} The self-similar solution has 4 parameters: $\alpha,m,d,M$. For each result and proof, the value of $d$ will not change, so we will not make this dependence explicit. When we intend to change the value of $\alpha,d,M$ through a proof statement, we will make this dependence explicit $U = U_{\alpha,m,M}$. Often we will drop the sub-indices in the interest of clarity. We will denote generic constants we will use the notation $c$ and $C$, and when their dependences can be inferred from the context we will.

\subsection{An integral equation for the self-similar solutions}

Computing the integral, we can show that the conservation of mass \eqref{eq:mass conservation} is equivalent to $a = db$ and
\begin{equation}
  \label{eq:profile mass}
  |\partial B_1| \int_0^\infty U(\rho) \rho^{d-1} d \rho  = M.
\end{equation}
We will show that the scaling factors of the equation are given by
\begin{equation}
  \label{eq:self-similar exponents}
  a = \frac{\alpha d}{2 + d(m-1)} , \qquad b = \frac{\alpha}{2+d(m-1)}.
\end{equation}
Notice that $a, b$ are positive if and only
\begin{equation}
  m>m_c \coloneqq (1-\tfrac{2}{d})_+.
\end{equation}
This restriction on $m$ will be kept in this paper.
Formally, if $u$ is classical solution of \eqref{eq:main} of the form \eqref{eq:ssf} then we will show that $U$ satisfies the integral operator equation
\begin{equation}
  \label{eq:profile equation}
  U(z)^{m} = \int_z^{+\infty} K(z,\rho) \, U(\rho) d\rho,
\end{equation}
with a non-negative kernel $ K(z,\rho)$ given for $\rho \ge z$ by
\begin{equation}
  \label{eq:Q}
  K(z,\rho)= \rho \, Q(\tfrac{z}{\rho}), \qquad Q(\eta) =
  \frac{1}{\Gamma(1-\alpha)}\int_\eta^1(1-\sigma^{\frac 1 b})^{-\alpha} \sigma^{1-d} d\sigma.
\end{equation}
Notice that $K(z,z) = 0$ and $K(z,\rho) > 0$ if $z < \rho$.
This formulation is a key tool in our paper.

\subsection{Existence of self-similar solutions for $m > m_c$}
Our first result is the existence of self-similar solutions and their characterization as weak solutions.
\begin{theorem}
  \label{thm:existence}
  Let $d \ge 1$, $\alpha \in (0,1)$, $m > m_c$, and $M > 0$. Then, there exists a canonical profile $\mathcal U_{\alpha,m} \in C((0,+\infty))$ such that $z^{d-1} \mathcal U_{\alpha,m}(z) \in L^1(0,\infty)$ and, if we let,
  %
  %
  %
  %
  \begin{equation}
    \label{eq:rescaling profile to correct mass}
    U_{\alpha,m,M}(z) = A \mathcal U_{\alpha,m}(A^{-\frac{m-1}{2}} z), \qquad  A(\alpha,m,M)^{1 + \frac{(m-1)d}{2}} = \frac{M}{|\partial B_1| \int_0^\infty \mathcal U_{\alpha,m}(\rho) \rho^{d-1} d \rho}.
  \end{equation}
  the following properties hold:
  \begin{enumerate}
    \item
      $\mathcal U_{\alpha,m}$ and $U_{\alpha,m,M}$ are solutions to \eqref{eq:profile equation} and $U_{\alpha,m,M}$ satisfies \eqref{eq:profile mass}.
    \item
      Letting $U = U_{\alpha,m,M}$, the self-similar solution $u = u_{\alpha,m,M}$ given by \eqref{eq:ssf}
      is a very weak solution of \eqref{eq:main} in the sense that
      \begin{equation}
        \label{eq:very weak solution}
        \begin{aligned}
          \int_{0}^T  \int_{\Rd} u(t,x) \dRL \varphi(T-t,x)dxdt
          & = \int_0^T \int_{\Rd} u^{m}(t,x) \Delta\varphi(T-t,x)  dx dt
          \\
          & \qquad + M I_{1-\alpha}\varphi(T,0).
        \end{aligned}
      \end{equation}
      for all $\varphi \in C^\infty_c ([0,T] \times \mathbb R^d)$. Furthermore, it satisfies \eqref{eq:initial Dirac} and \eqref{eq:profile mass}.
    \item
      \label{it:slow-diffusion U near 0}
      %
      Behaviour near $z = 0$: We have $z^{d-1}(U^m)' \in C([0,+\infty))$ and
      \begin{equation}
        \label{eq:limit of flux at z=0}
        - \lim_{z \to 0^+} |\partial B_1| z^{d-1}\frac{dU^m}{dz}(z) = \frac {M} {\Gamma(1-\alpha)}.
      \end{equation}
      Furthermore, we have that
      \begin{equation}
        \begin{aligned}
          \label{eq:profile at z=0}
          \lim_{z \to 0^+}\frac{U_{\alpha,m,M}^m(z)}{\mathcal V_{\alpha,m}^m(z)} =
          \begin{cases}
            \displaystyle\frac{M}{|\partial B_1| (d-2)\Gamma(1-\alpha)} & \text{if } d\ge 3,
            \vspace{2pt}\\
            \displaystyle\frac{M}{|\partial B_1| \Gamma(1-\alpha)} &\text{if } d = 2,
            \vspace{2pt}\\
            \displaystyle\frac{M \Gamma(b_{\alpha,m}+1)}{2 \Gamma(b_{\alpha,m}+1-\alpha)} &\text{if } d= 1.
          \end{cases} \text{ where } \mathcal V_{\alpha,m}^m(z) \coloneqq
          \begin{dcases}
            z^{-{d-2}}     & \text{if } d \ge 3 \\
            |\log {z}| & \text{if } d=2,     \\
            1                      & \text{if } d=1.
          \end{dcases}
        \end{aligned}
      \end{equation}
  \end{enumerate}
\end{theorem}

We divide the proof into cases $m = 1$ in \Cref{sec:existence proof m=1}, $m > 1$ in \Cref{sec:existence proof m>1}, and $m \in (m_c,1)$ in \Cref{sec:existence proof m < 1}.

\begin{remark}
  Notice that due to \eqref{eq:profile at z=0} we see that for $d \ge 2$ we have $\mathcal V_{\alpha,m}^m$ is the Newtonian potential up to a constant, and hence so is $u^m$.
  Therefore, $\Delta u^m(t,x) = \partial_t^\alpha u(t,x)$ is a multiple of $\delta_0$.
  This means that $\partial_t^\alpha$ always remembers the initial Dirac. 
  This was known for $m = 1$.
  In bounded domains and $0 < \alpha < 1$ it was shown in \cite{bonforteTimeFractionalPorousMedium2024} that if $u_0 \in L^p \cap H^*$ for $p \ge \max\{ \frac{d}{2s}, 1 + m\}$, then $u^m \in L^{\infty}$. The critical $p = \frac{d}{2s}$ is precisely the value such that Dirichlet Laplacian in bounded domains $-\Delta_\Omega$ satisfies $(-\Delta_\Omega)^{-1} : L^p \to L^\infty$. The authors had already pointed out that the memory effect makes
  $u^m$ behave like $(-\Delta_\Omega)^{-1} u_0$. This is \emph{elliptic regularization at work}.
  Lastly, we recall that the classical Barenblatt for $\alpha = 1$ is bounded and smooth at $z = 0$.
\end{remark}

\begin{remark}
  The theory of the existence of the profile in the linear case $m = 1$ is already well understood, and we only revise it in \Cref{sec:m=1}.
  We believe that \eqref{eq:very weak solution} and \eqref{eq:limit of flux at z=0} are also new for $m=1$.
\end{remark}

\subsection{The linear-diffusion case $m = 1$}
In \Cref{sec:m=1} we recall the deduction the closed formula
\begin{equation}
  \label{eq:m=1 self-similar}
  u_{\alpha,1,M}(t,x) = M \mathcal F^{-1}[ \widehat Z(t,\cdot) ] (x), \qquad \widehat Z(t,\xi) = E_\alpha(-|\xi|^2t^\alpha),
\end{equation}
where $E_\alpha$ is the Mittag-Leffler function. In dimension $d=1$, $Z$ can be written in terms of the Wright function, see \cite{MR1829592}.
Due to the explicit construction, it was proved in \cite{dier2020} for $\alpha \in (0,1)$
that for $z \ge 1 $ there exists constants such that
\begin{equation*}
  c \mathcal V_{\alpha,1} (z)
  \le
  {U_{\alpha,1,M}(z)}
  \le C\mathcal V_{\alpha,1}(z),
  \qquad
  \mathcal V_{\alpha,1}(z) \coloneqq z^{\frac{d(\alpha-1)}{2(2-\alpha)}} e^{-\sigma_\alpha z^{\frac 1{2-\alpha}}},
  \qquad
  \sigma_\alpha \coloneqq (2-\alpha)(\alpha^\alpha/4)^{1/(2-\alpha)}.
\end{equation*}
Further asymptotic analysis can be found in \cite{deng2019}.
It is known that the fundamental solution describes the intermediate asymptotics of the solution with general initial datum $u_0$, see, e.g. \cite{cortazarHeatEquationMemory2021}. The uniqueness of weak solutions holds easily by duality.
\begin{theorem}
  \label{thm:m=1 uniqueness}
  Let $d \ge 1, \alpha \in (0,1)$, $m = 1$, and $M > 0$. Then, there exists a unique solution to \eqref{eq:very weak solution}.
  Hence, there is at most one solution to \eqref{eq:profile equation} such that \eqref{eq:profile mass}.
\end{theorem}

\subsection{The slow-diffusion regime $m > 1$}
In the slow-diffusion case, we can a complete description of the self-similar profiles. Similarly to the case $\alpha = 1$ they are compactly supported and of power-type at the free boundary.
\begin{theorem}[Slow-diffusion regime]
  \label{thm:slow-diffusion}
  Let $d \ge 1$, $\alpha \in (0,1)$, and $m > 1$.
  Then, the profiles $\mathcal U_{\alpha,m}$ and $U_{\alpha,m,M}$ in \Cref{thm:existence} have the following properties
  \begin{enumerate}
    \item
      $U_{\alpha,m,M}$ is the unique non-negative point-wise solution to \eqref{eq:profile equation} such that \eqref{eq:profile mass}.
    \item The solutions are ordered according to their mass and for each $z > 0$
      \[
        U_{\alpha,m,M}(z) \to +\infty  \qquad \text{ everywhere as } M \to +\infty .
      \]
    \item
      \label{it:slow-diffusion U compactly supported}
      $\mathcal U_{\alpha,m}$ is continuous, non-increasing, and supported in $[0,1]$.
    \item
      %
      Behavior at the free boundary:
      \begin{equation}
        \label{eq:slow sharp constant at free boundary}
        \lim_{z \to 1^-} \frac{\mathcal U_{\alpha,m}(z)}{(1-z)_+^{\frac {2-\alpha }{m-1}}} = \left( \frac{b_{\alpha,m}^\alpha \Gamma(1 + \tfrac{2-\alpha}{m-1} )}{\Gamma(3-\alpha+\tfrac{2-\alpha}{m-1})} \right)^{\frac{1}{m-1}}.
      \end{equation}
    \item Classical limit as $\alpha \to 1^-$: for $z > 0$ we have
      \begin{equation}
        \label{eq:m>1 limit alpha->1}
        \lim_{\alpha \to 1^-}\mathcal U_{\alpha,m}(z) =  \left( \frac{(m-1)b_{1,m}}{2m} (1-z^2)_+ \right)^{\frac 1 {m-1}} \eqqcolon \mathcal U_{m}(z).
      \end{equation}
    \item
      \label{thm:mesa}
      The mesa limit: for any $t > 0$ and $x \in \mathbb R^d \setminus \{0\}$ we have that
      \begin{equation}
        \lim_{m \to \infty} u_{\alpha,m,M}(t,x) = M \frac{\chi_{\overline{B_1}} (x)}{|B_1|}  \eqqcolon u_{\infty,M}(x).
      \end{equation}
  \end{enumerate}
\end{theorem}

The proof of this result can be found in \Cref{sec:existence proof m>1}.
Notice that $\mathcal U_{m}$ is the famous Barenblatt solution; see, e.g. \cite[(17.28)-(17.30)]{Vazquez2006}.

\begin{remark}
  Notice that the mesa limit does not depend on $\alpha$. Moreover, it has lost the singularity at $x=0$ and does not change with time.
\end{remark}

\begin{figure}[h!]
  \centering
  \includegraphics[width=0.45\textwidth]{slow-mesa-d=1}
  \includegraphics[width=0.45\textwidth]{slow-mesa-d=3}
  \caption{Slow-diffusion profiles for different values of $m$}
  \label{fig:mesa}
\end{figure}

\begin{figure}[h!]
  \centering
  \includegraphics[width=0.45\textwidth]{slow-limit-alphas-d=1.pdf}
  \includegraphics[width=0.45\textwidth]{slow-limit-alphas-d=3.pdf}

  \caption{Slow-diffusion profiles for different values of $\alpha$.}
  \label{fig:limit alpha->1}
\end{figure}

%

\subsection{The mildly-fast-diffusion regime $m \in (m_c, 1)$}
We study the mildly-fast-diffusion regime $m_c<m<1$, where the critical value $m_c$ is the same as in the classical case $\alpha = 1$.
We recall that for $\alpha=1$ the self-similar profiles are $U_{1,m}(z) = c_{1,m}^* (z^2 + c)^{-\frac 1 {1-m}}$ where $c_{1,m}^*$ depends only on $m$ and $d$, and $c>0$ is a free constant.
These solutions have finite mass.
As $c \to 0^+$ we recover the so-called \emph{very-singular solution} (VSS), which has infinite mass. We will show that for $\alpha \in (0,1)$ there is the \emph{time-fractional very singular solution} given by
\begin{equation}\label{eq:vss}
  U_{\alpha,m}^*(z) = c_{\alpha,m}^* z^{-\frac 2 {1-m}}, \qquad  c_{\alpha,m}^* = \left( \frac{(1-m)b_{\alpha,m}}{2m} \frac{\Gamma(\tfrac{\alpha}{1-m})}{\Gamma(1-\alpha + \frac{\alpha}{1-m})}  \right)^{\frac{-1}{1-m}} > 0.
\end{equation}
This construction is justified in \Cref{lem:vss}. Notice that this function gives an exact self-similar solution in separated variables
\begin{equation*}
  u_{\alpha,m}^*(t,x) \coloneqq t^{-bd} U_{\alpha,m}^*(t^{-b} x) = c^*_{\alpha,m} \left(\frac{t^{{\alpha}}}{|x|^{2}}\right)^{\frac 1 {1-m}}.
\end{equation*}
This solution starts from a singularity point at $t=0$ and increases in time everywhere for all $x\ne 0$ for all $t>0$. 
The convention \emph{very-singular solution} is recent, and some previous literature has referred to this solution as \emph{infinite point-source solutions} (cf. \cite{Vazquez2006}).
%
%

%

%

%
%
%
%
%

%
%
%
We prove the existence of finite-mass self-similar solutions that form an \emph{asymptotic fan} around this very singular solution.

\begin{theorem}[Fast-diffusion regime, mildly-singular range]
  \label{thm:fast-diffusion}
  Assume that $d \ge 1$, $\alpha \in (0,1)$, $m \in (m_c, 1)$. Then the self-similar profiles from \Cref{thm:existence} are such that
  \begin{enumerate}
    \item $\mathcal U_{\alpha,m}$ and $U_{\alpha,m,M}$ are decreasing and strictly positive for $z \in (0,+\infty)$.

    \item The tail behavior is given by the asymptotic expansion
      \begin{equation}
        \label{eq:fast-diffusion U tail}
        \begin{aligned}
          \mathcal U_{\alpha,m}(z) & = U^*(z) ( 1 - z^{-\gamma^*} + \mathcal R_{\alpha,m}(z)),
          \\
          U_{\alpha,m,M}(z)        & = U^*(z) ( 1 - (Lz)^{-\gamma^*} + \mathcal R_{\alpha,m}(Lz))
        \end{aligned}
      \end{equation}
      where $L = A^{-\frac{m-1}{2}}$ is given by \eqref{eq:rescaling profile to correct mass}, $\mathcal R_{\alpha,m}(z) = o(z^{-\gamma^*})$, and the exact value of $\gamma^*$ is implicitly given by \eqref{eq:gamma^star}.
    \item The solutions are ordered with respect to $M$, and furthermore
      \begin{align*}
        & U_{\alpha,m,M}(z) \nearrow U_{\alpha,m}^*(z) \text{ as } M \to \infty
        \qquad \text{and} \qquad
        U_{\alpha,m,M}(z) \searrow 0 \text{ as } M \to 0.
      \end{align*}
    \item Classical limit as $\alpha \to 1^-$: We have that
      \begin{equation}
        \label{eq:m<1 alpha->1}
        \lim_{\alpha \to 1^-}\mathcal U_{\alpha,m} (z) = c^*_{1,m} (1+z^2)_+ ^{\frac 1 {m-1}} \eqqcolon \mathcal U_{m}(z).
      \end{equation}
  \end{enumerate}
\end{theorem}

\subsection{Continuity in $m$ at $m=1$.}
Our last result is the matching of the three behaviors. Due to the criteria by which we have chosen the ``canonical profile'' $\mathcal U_{\alpha,m}$ we do not expect continuity w.r.t. $m$ as $m$ crosses the critical value $m = 1$ (see \eqref{eq:m>1 limit alpha->1} and \eqref{eq:m<1 alpha->1}). 
Furthermore, it is easy to notice that $\mathcal U_{1,m} \to 0$ as $m \to 1^-$ and $\mathcal U_{1,m} \to +\infty$ as $m \to 1^+$.
However, the correct continuity in $m$ holds for the solutions of given mass, as shown by the following result.
\begin{theorem}
  \label{thm:limit m->1}
  Let $d \ge 1$, $\alpha \in (0,1)$, and $M > 0$. Then, for any $z > 0$ we have that
  \begin{equation*}
    \lim_{m \to 1} U_{\alpha,m,M}(z) = U_{\alpha,1,M}(z).
  \end{equation*}
\end{theorem}

%
\section{Weak formulation}

In this section we discuss the weak formulation, and provide a lemma that allows to characterize self-similar solutions as weak solution based on the properties of the profile.

\subsection{Weak formulation for smooth solutions}
We begin by showing that the weak formulation is reasonable for smooth solutions.
By Fubini, we deduce that
\begin{align*}
  \Gamma(\alpha) \int_\tau^T [I_{1-\alpha} f(t)] \, g(T-t) dt
  & = \int_\tau^T \int_0^t (t-s)^{-\alpha} f(s) g(T-t) ds dt
  \\
  & = \int_\tau^T \int_s^{T} (t-s)^{-\alpha} g(T-t) dt \, f(s) ds
  \\
  & = \int_\tau^T \int_0^{T-s} (T-\sigma -s)^{-\alpha} g(\sigma) d\sigma \, f(s)  ds
  \\
  & = \Gamma(\alpha) \int_\tau^T f(s) I_{1-\alpha} g (T-s) ds.
\end{align*}
This means that
\begin{equation*}
  \int_0^t \partial_t^\alpha f(t) g(T-t) dt = \int_0^T f(t) \dRL g (T-t)dt + f(T) I_{1-\alpha} g(0) - f(0) I_{1-\alpha} g(T).
\end{equation*}
Notice that if $g \in C([0,T])$ then $I_{1-\alpha} g(0) = 0$ for $\alpha \in (0,1)$.
For $\alpha = 1$ we have $I_{1-\alpha} g = g$.
We write therefore that
\begin{align*}
  \int_{0}^T \int_{\Rd}  u(t,x) \dRL  \varphi (T-t,x)dxdt
  &
  =
  \int_\Rd u(0,x) I_ {1-\alpha} \varphi(T,x) dx - \int_\Rd u(T,x) I_{1-\alpha}\varphi(0,x) dx \\
  & \quad +
  \int_0^{T} \int_{\Rd} u^{m} (t,x) \Delta \varphi (T- t,x) dxdt.
\end{align*}
Since $\varphi$ is bounded, we have $I_{1-\alpha}\varphi (0,x) = 0$ and therefore
\begin{equation*}
  \begin{aligned}
    \int_{0}^T \int_{\Rd}  u(t,x) \dRL  \varphi (T-t,x)dxdt
    &
    =
    \int_\Rd u(0,x) I_ {1-\alpha} \varphi(T,x) dx
    \\
    & \quad +\int_0^{T} \int_{\Rd} u^{m} (t,x) \Delta \varphi (T- t,x) dxdt.
  \end{aligned}
\end{equation*}
This integration by parts holds for all $\varphi \in C_c^\infty ([0,T] \times \mathbb R^d)$ since $\dRL$ does not notice finite values of $\varphi(0,x)$.

%
\subsection{Weak formulation for finite-mass self-similar solution}
Now we provide a general lemma to show that a self-similar solution is a weak solution in the sense of \eqref{eq:very weak solution}.
\begin{lemma}
  \label{lem:sss is weak solution}
  Let $d \ge 1$, $m > 0$, $\alpha \in (0,1)$, $u$ be given by \eqref{eq:ssf}.
  Assume that $U$ is supported in $[0,R]$ with $R \le \infty$ and satisfies \eqref{eq:limit of flux at z=0}, and that $u$ is a classical solution to \eqref{eq:main} in
  \begin{equation*}
    S \coloneqq \{(t,x) : t > 0, x \ne 0, u(t,x) > 0\},
  \end{equation*}
  such that \eqref{eq:mass conservation} for all $t > 0$, and that $\nabla u^{m}$ is continuous at $\partial \operatorname{supp}(u)$.
  Then $u$ satisfies \eqref{eq:very weak solution}.
\end{lemma}

\begin{proof}
  \textbf{First integration by parts.}
  Let us take $\varphi \in C_c^\infty([0,T] \times \mathbb R^d)$ and for $\delta \in (0,1)$ let us define
  \[
    \Omega_\delta(t) =
    \begin{cases}
      B_{(1-\delta)Rt^b} & \text{if } R < \infty , \\
      \mathbb R^d        & \text{if } R = \infty.
    \end{cases}
  \]
  Notice that if $x \in \Omega_\delta(t) \setminus B_\varepsilon $ and $t > 0$ then $(t,x) \in S$.
  Since $u$ is a classical solution in $S$, we have
  \begin{equation*}
    \int_{0}^T \int_{\Omega_\delta(t) \setminus B_\varepsilon} \partial_t^\alpha u(t,x)  \varphi (T-t,x)dxdt
    =
    \int_0^{T} \int_{\Omega_\delta(t) \setminus B_\varepsilon} \Delta u^{m} (t,x) \varphi (T- t,x) dxdt.
  \end{equation*}
  To integrate by parts the Caputo derivative, we write
  \begin{equation*}
    \int_{0}^T \int_{\Omega_\delta(t) \setminus B_\varepsilon} \partial_t^\alpha u(t,x)  \varphi (T-t,x)dxdt
    = \int_{\Omega_\delta(T) \setminus B_\varepsilon} \int_{A_\delta(x)}^T \partial_t^\alpha u(t,x)  \varphi (T-t,x)dtdx.
  \end{equation*}
  We repeat the integration by parts of the Caputo derivative
  \begin{align*}
    \int_A^T \partial_t^\alpha f(t) g(T-t) dt & = \int_0^T \partial_t f(t) I_{1-\alpha} g (T-t) dt - \int_0^A \partial_t f (t) R(t) dt,
  \end{align*}
  where
  \[
    R(t) \coloneqq \int_0^A \int_{T-A}^{T-t} \frac{(T-t-\sigma)^{-\alpha}}{\Gamma(1-\alpha)} \varphi(\sigma) d \sigma.
  \]
  Integrating by parts, we recover
  \begin{align*}
    \int_A^T \partial_t^\alpha f(t) g(T-t) dt
    & = \int_0^T  f(t) \dRL g (T-t) dt  - f(0) I_{1-\alpha} g(T)          \\
    & \qquad - \int_0^A f (t) \partial_t R(t) dt + f(A) R(A) - f(0) R(0).
  \end{align*}
  Since $u(0,x)= 0$ in $\Rd \setminus B_\varepsilon$ we obtain as $\delta \to 0$ that
  \begin{align*}
    \int_{\tau}^T \int_{\Omega_\delta(t) \setminus B_\varepsilon} \partial_t^\alpha u(t,x)  \varphi (T-t,x)dxdt
    & \overset{\delta \to 0}\longrightarrow
    \int_{\tau}^T \int_{B_{Rt^\beta} \setminus B_\varepsilon} u(t,x) \dRL \varphi (T-t,x)dxdt
    \\
    & =
    \int_{\tau}^T \int_{\Rd \setminus B_\varepsilon} u(t,x) \dRL  \varphi (T-t,x)dxdt.
  \end{align*}
  On the other hand, we have that
  \begin{align*}
    \int_0^{T} \int_{\Omega_\delta(t) \setminus B_\varepsilon} \Delta u^{m} (t,x) \varphi (T- t,x) dxdt
    & = - \int_0^{T} \int_{\Omega_\delta(t) \setminus B_\varepsilon} \nabla u^{m} (t,x) \nabla \varphi (T- t,x) dxdt
    \\
    & \qquad + \int_0^T \int_{\partial B_\varepsilon} \varphi(T-t,x) \nabla u^{m}(t,x) \cdot {\frac{-x}{|x|}}
    \\
    & \qquad + \int_0^T \int_{\partial B_{(1-\delta)Rt^b}} \varphi(T-t,x) \nabla u^{m}(t,x) \cdot {\frac{x}{|x|}}.
  \end{align*}
  Since $\nabla u^{m}$ is continuous at $\partial \operatorname{supp}(u)$ and $u(t,x) = 0$ if $x \notin \Omega_\delta(t)$ then we have that
  \begin{align*}
    \int_0^{T} \int_{\Omega_\delta(t) \setminus B_\varepsilon} \Delta u^{m} (t,x) \varphi (T- t,x) dxdt
    & \overset{\delta \to 0} \longrightarrow
    - \int_0^{T} \int_{\mathbb R^d \setminus B_\varepsilon} \nabla u^{m} (t,x) \nabla \varphi (T- t,x) dxdt
    \\
    & \qquad + \int_0^T \int_{\partial B_\varepsilon} \varphi(T-t,x) \nabla u^{m}(t,x) \cdot {\frac{-x}{|x|}}.
  \end{align*}
  We conclude that
  \begin{equation}
    \label{eq:weak formulation 1}
    \begin{aligned}
      \int_{0}^T \int_{\Rd \setminus B_\varepsilon} u(t,x) \dRL \varphi (T-t,x)dxdt
      & =
      - \int_0^{T} \int_{\mathbb R^d \setminus B_\varepsilon} \nabla u^{m} (t,x) \nabla \varphi (T- t,x) dxdt
      \\
      & \qquad
      + \int_0^T \int_{\partial B_\varepsilon} \varphi(T-t,x) \nabla u^{m}(t,x) \cdot {\frac{-x}{|x|}}.
    \end{aligned}
  \end{equation}
  Notice that the last term vanishes as $\tau \to 0$ for $\varepsilon$ fixed.
  %

  \textbf{Limit as $\varepsilon \to 0$.} %
  Since $u \in L^\infty(0,T; L^1 (\mathbb R^{d}))$ we have that
  \[
    \int_{0}^T \int_{\Rd \setminus B_\varepsilon} u(t,x) \dRL \varphi (T-t,x)dxdt
    \overset{\varepsilon \to 0} \longrightarrow
    \int_{0}^T \int_{\Rd} u(t,x) \dRL \varphi (T-t,x)dxdt.
  \]
  We notice that $\nabla u^m$ is not a nice function in $x = 0$. We integrate by parts again in space
  \begin{align*}
    \int_0^T \int_{\Rd \setminus B_\varepsilon} u(t,x) \dRL \varphi (T-t,x)dxdt
    & =
    \int_0^{T} \int_{\mathbb R^d \setminus B_\varepsilon} u^{m} (t,x) \Delta \varphi (T- t,x) dxdt
    \\
    & \qquad
    + \int_{\Rd \setminus B_\varepsilon} u(\tau,x) I_{1-\alpha} \varphi (T,x)dxdt.
    \\
    & \qquad
    + \int_0 ^T \int_{\partial B_\varepsilon} \varphi(T-t,x) \nabla u^{m}(t,x) \cdot {\frac{-x}{|x|}}
    \\
    & \qquad
    -
    \int_0^T \int_{\partial B_\varepsilon} u^{m}(t,x) \nabla \varphi(T-t,x) \cdot {\frac{-x}{|x|}}.
  \end{align*}
  Since $|\partial B_\varepsilon| \le C \varepsilon^{d-1}$, the last term is bounded by $\varepsilon \|\nabla \varphi\|_{L^\infty}$.
  For the second-to-last term, we again use \eqref{eq:limit of flux at z=0} in its weaker form.
  Again we write that
  \begin{align*}
    \int_{\partial B_\varepsilon} \varphi(T-t,x) \nabla u^{m}(t,x) \cdot \frac{-x}{|x|}
    & = G(t^{-b}\varepsilon) t^{-\alpha} |B_\varepsilon|^{-1} \int_{\partial B_\varepsilon} \varphi(T-t, x) dx
    ,
  \end{align*}
  where
  $G(z) = \frac{dU^{m}}{dz} z^{d-1}$.
  Now we use the smoothness of $\varphi$. Let $\eta(t) = \varphi(t,0)$ and we have that
  \begin{equation*}
    \sup_{t\in [0,T]} \left| \eta(t) - |\partial B_\varepsilon|^{-1} \int_{\partial B_\varepsilon} \varphi(t,x) \right| \le \|\nabla \varphi\|_{L^\infty} \varepsilon.
  \end{equation*}
  Hence, we have that
  \begin{equation*}
    \left| \int_0^T \int_{\partial B_\varepsilon} \varphi(T-t,x) \nabla u^{m}(t,x) \cdot \frac{-x}{|x|} - \int_0^T \eta(T-t) t^{-\alpha} G(t^{-b} \varepsilon) dt \right|  \le C \varepsilon.
  \end{equation*}
  As $\varepsilon \to 0$ we use \eqref{eq:limit of flux at z=0} to deduce that
  \[
    \int_0^T \int_{\partial B_\varepsilon} \varphi(T-t,x) \nabla u^{m}(t,x) \cdot \frac{-x}{|x|} \overset{\varepsilon \to 0} \longrightarrow \frac{M}{\Gamma(1-\alpha)} I_{1-\alpha} \varphi(T,0).
  \]
  In the end, we conclude \eqref{eq:very weak solution}.
\end{proof}

%
\section{Equation for the self-similar profile}
In this section, we will derive the integral equation that represents the self-similar solution and prove several of its properties.

%

\subsection{Scaling properties of \eqref{eq:main}}
\label{sec:scaling}
First, let us obtain the scaling factors $a$ and $b$ through scaling arguments at the PDE level.
This equation has the following scaling. If $u(t,x)$ is a solution of a given mass $M$, then
\begin{equation*}
  \widetilde u(t,x) = A u(Tt, Lx ),
\end{equation*}
is also a solution if the three non-negative scaling constants are related by
\begin{equation}
  \label{eq:dimensional scaling}
  A^{m-1} L^2 = T^\alpha.
\end{equation}
Hence, if $U$ is a profile, then
\[
  \widetilde U (z) = A T^{-bd} U(T^{-b} L z),
\]
is also a self-similar profile of the same equation. We point out that
\[
  \int_\Rd \widetilde u(t,x) dx = A L^{-d} \int_\Rd u(Tt, y) dy,
\]
so the scaling of the masses is
\[
  \widetilde M = A L^{-d} M.
\]
Hence, once we show the existence of a profile of a certain mass, we will immediately have the existence of profiles of all masses.

To obtain a self-similar scaling, we set $\widetilde U = U$. We must therefore have $A L^{-d} = AT^{-bd} = T^{-b}L = 1$ so $A = T^{bd}$ and $L = T^b$. Lastly, matching the scaling of equation \eqref{eq:dimensional scaling}, we deduce that $bd(m-1) + 2b = \alpha$, hence the correct scaling of dimensions \eqref{eq:self-similar exponents}. On the other hand, if we aim to change mass, we set $T = 1$, and we have that
\begin{equation}
  \label{eq:scaling solutions}
  \widetilde U (z) = A U(A^{-({m-1})/2} z), \qquad L = A^{-({m-1})/2}
\end{equation}
is also a solution for any $A > 0$.
We are ready to discuss the existence of a suitable profile $U\ge 0$ with finite mass and compact support.

\subsection{Derivation of \eqref{eq:profile equation}}
Let us obtain the equation for the profile by inserting formula \eqref{eq:ssf} into equation \eqref{eq:main}. 
We follow the same approach as \cite{CaballeroOkrasinska-PlociniczakPlociniczakSadarangani2025}, but now in general dimension $d \ge 1$ and keeping track that reasoning is valid also for $m \in (m_c,1]$.

\paragraph{Scaling the Caputo derivative.} 
For $x \ne 0$ we have $u(0,x) = 0$ hence we can move the derivative outside the integral and write \eqref{eq:Caputo}
\begin{align*}
  \partial_t^\alpha u(t,x)
  & = I_{1-\alpha} \partial_t u (t,x)                                                                                                           = \partial_t I_{1-\alpha} u (t,x) %
  %
  %
  %
  %
  %
  %
  %
  %
  %
  %
  %
  = t^{-\alpha-a} \left( \mathsf A F(z) - b z \frac{dF}{dz}  \right)\Bigg|_{z = t^{-b}|x|}
\end{align*}
where $\mathsf A = 1-\alpha-a$ and
\[
  F(z) = \frac 1  {\Gamma(1-\alpha)}\int_0^1 (1-\sigma)^{-\alpha} \sigma^{-a} U(\sigma^{-b} z) d\sigma.
\]
Notice that, even though the integrand contains a singular factor $\sigma^{-a}$ with possibly $a \ge 1$,  this integral is well-defined if $|x| > 0$ and $U$ decays sufficiently fast at infinity.

\paragraph{Scaling the diffusion.} The spatial part, on the other hand, under the scaling behaves as follows
\[
  \Delta u^{m} (t,x) = t^{-2b-am} \left[ z^{1-d} \frac{d}{dz}\left(z^{d-1} \frac{dU^{m}}{dz}\right) \right] \Bigg|_{z=t^{-b} |x|}.
\]

\paragraph{An equation in self-similar variable.}
At this stage, the self-similar exponents \eqref{eq:self-similar exponents} can be deduced again.
Collecting the above scaling relations, we arrive at the integro-differential equation
\begin{equation}
  \label{eq:profile equation second order}
  z^{1-d} \frac{d}{dz} \left(z^{d-1} \frac{dU^{m}}{dz} \right) = \mathsf A F(z) - b z \frac{dF}{dz}.
\end{equation}
At this stage, this is a formal derivation. For the profiles we construct, we will show that this equation holds pointwise in the interior of the support and we will show that $u$ is a weak solution of \eqref{eq:main}.

\begin{comment}
\todo[inline]{
    \color{teal}
    Should we have $m+1$ or $m$ in the derivative? If $U$ has a compact support, then $dU^{m+1}/dz$ and $F$ vanish at $\infty$ from definition. Then, these assumptions are redundant. Is that correct? Or, maybe we want to include the fast-diffusion tail assumption here $z^d F(z) \rightarrow \infty$?
    \normalcolor
    \color{blue}
    It should be $U^m$.
    We should unify the presentation for the fast and slow diffusion. {\color{teal}Agreed!} JLV OK
    \normalcolor
}
\end{comment}

\paragraph{An integral equation.} If we have that
\begin{equation}
  \lim_{z \to \infty} z^{d-1}\frac{dU^{m}}{dz}(z) = 0,
  \qquad
  \lim_{z\to \infty} z^d F(z) = 0,
\end{equation}
then we have
\begin{align*}
  -z^{d-1} \frac{dU^{m}}{dz} (z) = \mathsf A \int_z^{+\infty} F(\sigma) \sigma^{d-1}d\sigma - b \int_{z}^{+\infty} \sigma^d \frac{dF}{dz}(\sigma) d \sigma.
\end{align*}
Integrating by parts and noticing that $\mathsf A + db = 1-\alpha$ we have that
\begin{equation}
  \label{eq:equation for dUm/dz}
  -z^{d-1} \frac{dU^{m}}{dz} (z) =  \int_z^{+\infty} ({1-\alpha}) F(\sigma) \sigma^{d-1} d\sigma + b z^{d} F(z).
\end{equation}

\begin{remark}
  \label{rem:decreasing}
  If a solution of \eqref{eq:equation for dUm/dz}
  is $U \ge 0$,  then $U$ is non-increasing.
\end{remark}
Lastly assuming that $U(+\infty) = 0$ we have that
\begin{equation}
  \label{eq:implicit equation last formula with F}
  \begin{aligned}
    U(z)^{m} & = \int_z^{+\infty} \sigma^{1-d} \int_\sigma^R ({1-\alpha}) \xi^{d-1} F(\xi) d \xi d \sigma + b \int_{z}^R \xi F(\xi)d\xi \\
    & = \int_z^{+\infty} \Bigg( ({1-\alpha}) \int_z^\xi \sigma^{1-d} \xi^{d-1} d \sigma + b \xi \Bigg) F(\xi) d \xi            \\
    & = \int_z^{+\infty} \Bigg( ({1-\alpha}) \xi \int_{z/\xi}^1 \eta^{1-d} d \eta + b \xi \Bigg) F(\xi) d \xi                  \\
    & = \int_z^{+\infty} \Bigg(({1-\alpha}) H_d(z/\xi) + b \Bigg)\xi F(\xi) d \xi,
  \end{aligned}
\end{equation}
where for $s \in (0,1)$ we have the Newtonian potentials (up to a constant)
\[
  H_d(s) = \int_{s}^1 \tau^{1-d} d\tau =
  \begin{cases}
    1-s                   & \text{if } d=1,     \\
    \log (1/s)            & \text{if } d = 2,   \\
    \frac{s^{2-d}-1}{d-2} & \text{if } d \ge 3.
  \end{cases}
\]
\begin{remark}
  Notice that in this last integration step it is not easy to integrate in $[0,z]$ and write $-(U(0)^{m} - U(z)^{m})$, because instead of $H_d$ we would have $\int_0^s \tau^{1-d}$.
\end{remark}
Taking into account that the solution has a compact support $[0,R]$, the lower limit of the integral is actually larger than $0$, therefore
\[
  U(z)^{m} = \frac{1}{\Gamma(1-\alpha)}\int_z^{+\infty} \left(({1-\alpha}) H_d(z/\xi) + b  \right) \xi\left( \int_{0}^1 (1-\sigma)^{-\alpha} \sigma^{-a} U(\sigma^{-b} \xi) d \sigma \right) d \xi.
\]
By substituting $\rho = \sigma^{-b}\xi$ and using the compact support, we obtain
\[
  \begin{split}
    U(z)^{m}
    & = \frac{1}{b}\frac{1}{\Gamma(1-\alpha)}\int_z^R \left(({1-\alpha}) H_d(\tfrac z \xi) + b \right) \xi^{1 + \frac{1-a}{b}} \left( \int_\xi^R \left(1-\left(\tfrac{\xi}{\rho}\right)^\frac{1}{b}\right)^{-\alpha} \rho^{\frac{a-1}{b}-1} U(\rho)d\rho\right) d \xi                        \\
    & = \frac{1}{b}\frac{1}{\Gamma(1-\alpha)}\int_z^R \left(\int_z^\rho \left(({1-\alpha}) H_d\left(\tfrac{z}{\xi}\right) + b \right) \left(\tfrac{\xi}{\rho}\right)^{1 + \frac{1-a}{b}} \left(1-\left(\tfrac{\xi}{\rho}\right)^\frac{1}{b}\right)^{-\alpha}d\xi \right)  U(\rho)d\rho.
  \end{split}
\]
%
%
Now, recalling that $\mathsf A+d b = (1-\alpha)$ we can write
\begin{equation*}
  U(z)^{m} = \int_z^R K(z,\rho)\, U(\rho) \,d\rho,
\end{equation*}
where the kernel $K(z,\rho)\ge 0$ acts only for $0<z<\rho$ and is given by the exact formula:
\begin{align*}
  K(z,\rho) & := \rho\, Q\left(\tfrac z \rho\right),
  \quad
  Q(\eta) =
  \frac{1}{\Gamma(1-\alpha)} \int_{\eta}^1 \left(\frac{1-\alpha}{b} H_d\left(\eta \sigma^{-1} \right) + 1 \right) \sigma^{1 + \frac{1}{b}-d} \left(1-\sigma^\frac{1}{b}\right)^{-\alpha} d\sigma.
\end{align*}
We arrive at the formulation \eqref{eq:profile equation}.
This is the integral equation that will serve as the basis for the study of the existence, uniqueness, and main properties of $U$.
The equation holds for all $m > 0$.

\paragraph{Simplification to \eqref{eq:Q}.}
Since $0<\eta<1$ we get
\begin{align*}
  \Gamma(1-\alpha) \,  Q'(\eta)
  & = -\left(\frac{1-\alpha}{b} H_d\left(1 \right) + 1 \right) \eta^{1 + \frac{1}{b}-d} \left(1-\eta^\frac{1}{b}\right)^{-\alpha}
  \\
  & \qquad +\int_{\eta}^1 \frac{1-\alpha}{b} H_d'\left(\eta \sigma^{-1} \right) \sigma^{\frac{1}{b}-d} \left(1-\sigma^\frac{1}{b}\right)^{-\alpha} d\sigma.
\end{align*}
Since $H_d(1) = 0$ and $H_d'(\sigma) = -\sigma^{1-d}$  we arrive at
\begin{equation*}
  K_z(z,\rho)=
  Q'(\eta) =
  -\frac{1}{\Gamma(1-\alpha)} \eta^{1 + \frac{1}{b}-d} \left(1-\eta^\frac{1}{b}\right)^{-\alpha}
  -\eta^{1-d}\frac{1-\alpha}{b\Gamma(1-\alpha)} \, \int_{\eta}^1 \sigma^{ \frac{1}{b}-1} \left(1-\sigma^\frac{1}{b}\right)^{-\alpha} d\sigma.
\end{equation*}
Hence $Q'(\eta) < 0$ for $0<\eta<1$, and we have $K_z(z,\rho)<0$ for $0<z<\rho$. Note that when $d=1$ the integral of $Q'(\eta)$ simplifies since $H'_1(s)=-1$ for $0<s<1$.
Furthermore, we can integrate
$$
\int_{\eta}^1   \sigma^{ \frac{1}{b}-1} \left(1-\sigma^\frac{1}{b}\right)^{-\alpha} d\sigma=
b \int_{\eta^{1/b}}^1   t^{1-b} (1-t)^{-\alpha} t^{b-1} dt = \frac{b}{1-\alpha}(1-\eta^{1/b})^{1-\alpha}.
$$
Simplifying we deduce that
\begin{equation}
  \label{eq:Q'}
  -Q'(\eta) =
  \frac{1}{\Gamma(1-\alpha)}\,(1-\eta^{\frac 1 b})^{-\alpha} \eta^{1-d}.
\end{equation}
Furthermore, since $Q(1) = 0$ we have \eqref{eq:Q}.

\begin{remark}
  We can introduce the incomplete beta notation for $x \in (0,1)$
  \begin{equation*}
    B_x(p,q) = \int_0^x \sigma^{p-1} (1-\sigma)^{q-1} d \sigma, \qquad I_x(p,q) = \frac {B_x(p,q)} {B(p,q)} .
  \end{equation*}
  Even though for $x < 1$ the value of $q$ is irrelevant, the literature only discusses these functions for $p,q > 0$.Using the change of variable $\sigma = \xi^b$ and $\eta=1-\xi$ we get
  \begin{equation*}
    Q(\eta) = \int_{\eta}^1 Q'(\sigma) d \sigma = \frac b {\Gamma(1-\alpha)} \int_{\eta^{1/b}}^1 {(1-\xi)^{-\alpha} \xi^{b(2-d) - 1 }} d \xi = \frac b {\Gamma(1-\alpha)} \int_0^{1-\eta^{1/b}} {\sigma^{-\alpha} (1-\sigma)^{b(2-d) - 1 }} d \sigma.
  \end{equation*}
  We conclude that in dimension $d=1$ we can write
  \begin{equation}
    \label{eq:Q from incomplete gammas}
    Q(\eta) = \frac{b}{\Gamma(1-\alpha)} B_{1-\eta^{1/b}}\Big(1-\alpha, b\Big) = \frac{b \Gamma(b)}{\Gamma(1-\alpha + b)} I_{1-\eta^{1/b}}\Big(1-\alpha, b\Big).
  \end{equation}
\end{remark}

\paragraph{Asymptotic of $Q$ and $\eta = 0, 1$.}
Lastly, we point out the limits
For all dimensions $d \ge 1$, the singularity is dominated by the fractional boundary behavior, yielding the optimal constant
\[
  \lim_{\eta \to 1^-} \frac{Q_{\alpha,m}(\eta)}{(1-\eta)^{1-\alpha}} = \frac{b^\alpha}{\Gamma(2-\alpha)}.
\]
%
%
%
%
%
%
The singular behaviour at $\eta = 0$ is given by the power of the reference profile
\begin{equation}
  \label{eq:Q at 0}
  \lim_{\eta \to 0^+} \frac{Q_{\alpha,m}(\eta)}{\mathcal V_{\alpha,m}(\eta)^m} =
  \begin{dcases}
    \frac{1}{(d-2)\Gamma(1-\alpha)} & \text{if } d \ge 3 , \\
    \frac{1}{\Gamma(1-\alpha)} & \text{if } d =2 , \\
    \frac{\Gamma(b+1)}{\Gamma(b+1-\alpha)} & \text{if } d=1.
  \end{dcases}
\end{equation}

%
%
%
\subsection{Integro-differential version. Monotonicity and equi-continuity}
Let us move to the study of the properties of the solution to \eqref{eq:profile equation}. If we differentiate with respect to $z$ we get the equivalent formulation
\begin{equation}
  \label{intdiff. equation}
  (U(z)^{m})' = \int_z^{+\infty} K_z(z,\rho) \, U(\rho) d\rho,
\end{equation}
Recalling the above formulas, we get
$$
K_z(z,\rho)= \rho \, Q'(\eta)\frac1{\rho}= Q'(\eta), \quad \text{ where } \eta := \frac{z}{\rho}.
$$
Since $Q' < 0$ we deduce that
\begin{corollary}
  Any $U \in C((0,\infty))$ solution of \eqref{eq:profile equation} is strictly decreasing in its support.
\end{corollary}

Since $U$ is monotone, we can apply Lieb's trick
\begin{equation*}
  M = |\partial B_1| \int_0^z U(\rho) \rho^{d-1} d \rho  \ge U(z) |\partial B_1| \frac{z^d}{d}.
\end{equation*}
Hence, we recover the universal local $L^\infty$ bound depending only on the mass
\begin{equation}
  \label{eq:Lieb upper bound}
  U_{\alpha,m,M} (z) \le z^{-d} \frac{dM}{|\partial B_1|} \qquad \forall z > 0.
\end{equation}
From this estimate we also deduce the equi-continuity.
We point out that
\begin{equation*}
  \left| z^{d-1} \frac{dU^m}{dz} (z) \right| = \int_z^\infty (-Q')(\tfrac{z}{\rho}) \frac{z^{d-1}}{\rho^{d-1}} \rho^{d-1} U(\rho) d \rho \le \frac{1}{\Gamma(1-\alpha)}\int_{z}^\infty U(\rho)\rho^{d-1}d\rho .
\end{equation*}
Therefore,
\begin{equation}
  \label{eq:equicontinuity}
  \left| z^{d-1} \frac{dU_{\alpha,m,M}^m}{dz} (z) \right| \le \frac{M}{|\partial B_1| \Gamma(1-\alpha)}.
\end{equation}

%
%
%
%

%
%
%
%
%
%
%
%
%
%

%
%
%
%
%
%
%
%

%
\subsection{Analysis of the singularity of $(U^m)'$ at $z = 0$. Proof of \Cref{eq:limit of flux at z=0}}

The aim of this is to justify \eqref{eq:limit of flux at z=0}.
First, we provide a formal argument through the PDE.
Then, we obtain a direct proof for solutions of \eqref{eq:profile equation} using the properties of $Q$.

\paragraph{Formal estimate through the PDE.}
Let $u$ be given by \eqref{eq:ssf}. If we consider $\varphi(t) = t^{\alpha - 1}$ and multiply the equation by $\varphi(T-t)$
\begin{align*}
  \int_\tau^T \int_{\Rd \setminus B_\varepsilon} (T-t)^{\alpha - 1} I_{1-\alpha} \partial_t u(t,x) dx dt = \int_0^T (T-t)^{\alpha - 1} \int_{\Rd \setminus B_\varepsilon} \Delta u^{m} dx dt.
\end{align*}
Integrating by parts, we recover that
\begin{align*}
  \int_\tau^T  (I_{1-\alpha}\varphi)(T-t)  \partial_t \int_{\Rd \setminus B_\varepsilon} u(t,x) dx dt
  & = \int_\tau^T (T-t)^{\alpha - 1} \int_{\partial B_\varepsilon} \nabla u^{m} \cdot \frac{-x}{|x|}.
  %
  %
\end{align*}
By scaling the integrals and using the properties of Euler's beta function we deduce that $I_{1-\alpha} \varphi = \Gamma(\alpha)$.
Therefore, we observe that
\begin{align*}
  \Gamma(\alpha)
  & \left( \int_{\Rd \setminus B_\varepsilon} u(t, x)dx - \int_{\Rd \setminus B_\varepsilon} u(\tau,x)dx  \right)
  \\
  & = -\int_\tau^T (T-t)^{\alpha - 1} |\partial B_1| \varepsilon^{d-1} t^{-bd(m+1)}\frac{dU^{m}}{dz}(t^{-b}\varepsilon) t^{-b} dt
  \\ &\approx \left(-|\partial B_1| z^{d-1} \frac{dU^{m}}{dz} \right)\Bigg|_{z = 0^+} \int_\tau^T (T-t)^{\alpha - 1} t^{-\alpha} dt.
  %
  %
\end{align*}
We can now let $\tau \to 0$, where $u(\tau,x) = 0$ in $\Rd \setminus B_\varepsilon$ for $\tau$ small enough, and we can lastly let $\varepsilon \to 0$ to recover \eqref{eq:limit of flux at z=0}.

%

\paragraph{Through the self-similar equation.}
We start with the following result about the behavior of the derivative at the origin.
\begin{lemma}
  \label{lem:limit of flux holds}
  Let $U \in C((0,+\infty))$ be a finite mass solution.
  Then, $z^{d-1} (U^m)' \in C([0,+\infty))$ and \eqref{eq:limit of flux at z=0}.
\end{lemma}
\begin{proof}
  Let us denote, for this proof, $G(\eta) = \eta^{d-1} Q'(\eta) = \frac{(1-\eta^{\frac 1 b})^{-\alpha}}{\Gamma(1-\alpha)}$ which is non-decreasing in $\eta$.
  For $z > 0$ we write
  \[
    z^{d-1} (U^m)' =
    \int_z^\infty G(\tfrac z\rho) \rho^{d-1} U(\rho) d\rho
  \]
  since $G$ is continuous and $U(\rho) \rho^{d-1}$ is continuous and integrable, we recover the continuity of $z^{d-1} (U^m)'$.
  In order to pass to the limit, we write
  \begin{equation}
    \label{eq:limit of flux holds decomposition}
    \int_z^\infty G(\tfrac z\rho) \rho^{d-1} U(\rho) d\rho =
    G(0^+) \int_{z}^{+\infty} U(\rho) \rho^{d-1} d \rho + \int_z^\infty (G(\tfrac{z}{\rho}) - G(0^+)) U(\rho) \rho^{d-1}d\rho.
  \end{equation}
  The first term on the right-hand side has a limit because $U$ is non-negative.
  %
  %
  %
  %
  %
  %
  %
  %
  %
  %
  To show that the second term vanishes, we perform the change of variables $\eta = z/\rho$, which yields $d\rho = -z/\eta^2 d\eta$:
  \begin{align*}
    \int_z^\infty (G(\tfrac{z}{\rho}) - G(0^+)) U(\rho) \rho^{d-1}d\rho
    &= \int_0^1 (G(\eta) - G(0^+)) U(z/\eta) (z/\eta)^{d-1} \frac{z}{\eta^2} d\eta \\
    &= \int_0^1 \frac{G(\eta) - G(0^+)}{\eta} \left[ U(z/\eta) (z/\eta)^d \right] d\eta.
  \end{align*}
  Let us define $\Phi(\eta) = \frac{G(\eta) - G(0^+)}{\eta}$ and $F(\rho) = U(\rho)\rho^d$.

  We first analyze the integrability of $\Phi(\eta)$ on $(0,1)$. As $\eta \to 1^-$, we have $\Phi(\eta) \sim (1-\eta)^{-\alpha}$, which is locally integrable since $\alpha \in (0,1)$. As $\eta \to 0^+$, a Taylor expansion yields $G(\eta) - G(0^+) \sim \eta^{1/b}$, and thus $\Phi(\eta) \sim \eta^{1/b - 1}$. Because we are in the range where $b > 0$, this singularity is also locally integrable. Hence, $\Phi \in L^1(0,1)$.

  Next, we consider $F(\rho)$. Because $U$ is non-increasing and $U(\rho)\rho^{d-1} \in L^1(0,\infty)$, we can bound $F(\rho)$ by observing that
  \begin{equation*}
    F(\rho) = U(\rho) \rho^d \le d \int_0^\rho U(s) s^{d-1} ds.
  \end{equation*}
  We conclude that $F(\rho)$ is globally bounded, $F(\rho) \le C(d) M$, and that $\lim_{\rho \to 0^+} F(\rho) = 0$.

  Therefore, for any fixed $\eta \in (0,1)$, the integrand satisfies $\Phi(\eta) F(z/\eta) \to 0$ pointwise as $z \to 0^+$. Furthermore, the integrand is uniformly bounded by $M \Phi(\eta) \in L^1(0,1)$. By the dominated convergence theorem, the second term in the right-hand side of \eqref{eq:limit of flux holds decomposition}  vanishes in the limit $z \to 0^+$.
\end{proof}
\begin{comment}
This a priori estimate shows the behavior of the profile at $z = 0$
\begin{enumerate}
  \item In dimension $d = 1$ the profile is bounded, but has a cusp at $z = 0$
  \item In dimension $d > 1$ the profile is singular at $z = 0$. The derivative always behaves as $dU^{m}/dz \sim -z^{1-d}$. Therefore, we have the following
    \begin{equation*}
      U \sim
      \begin{dcases}
        (-\log z) ^{\frac{1}{m}} & \text{if } d = 2,   \\
        z^{-\frac{d-2}{m}}       & \text{if } d \ge 3,
      \end{dcases} \qquad \text{for } z \sim 0^+.
    \end{equation*}
\end{enumerate}
When $\alpha = 1$ this derivative is $0$ which confirms the well-known result that the classical Barenblatt solutions are smooth at $x = 0$.
\end{comment}
%
%
%
%
%
%
%
%

%

\subsection{Analysis of the singularity of $U$ at $z = 0$. Proof of \Cref{eq:profile at z=0}}
\label{sec:U at z=0}
We now use \eqref{eq:Q at 0} to deduce \eqref{eq:profile at z=0}, similarly to \Cref{lem:limit of flux holds}. Since the argument is analogous, we will not write a detailed proof.
\paragraph{Case $d \ge 3$.} We have that
\begin{align*}
  \frac{U_{\alpha,m,M}(z)^m}{\mathcal V_{\alpha,m}(z)^m} &= \frac{ U (z)^m }{z^{2-d}} = \int_{z}^\infty (\tfrac{z}{\rho} )^{d-2}Q(\tfrac{z}{\rho})  U(\rho) \rho^{d-1} d \rho
  \\
  &\to \frac{M}{|\partial B_1| (d-2)\Gamma(1-\alpha)} .
\end{align*}
As $\alpha \to 1^-$ this limit vanishes because $U_{1,m,M}(z)$ is bounded at $z = 0$.

\paragraph{Case $d =2$.} We have that
\begin{align*}
  \frac{U_{\alpha,m,M}(z)^m}{\mathcal V_{\alpha,m}(z)^m} &= \frac{ U (z)^m }{-\log z}
  = \int_{z}^\infty \frac{1}{-\log \frac z \rho - \log \rho} Q(\tfrac{z}{\rho})  U(\rho) \rho^{d-1} d \rho
  \\
  &\to \frac{M}{|\partial B_1| \Gamma(1-\alpha)} .
\end{align*}
As $\alpha \to 1^-$ this limit vanishes because $U_{1,m,M}(z)$ is bounded at $z = 0$.

\paragraph{Case $d = 1$.} We have that
\begin{equation*}
  \frac{U_{\alpha,m,M}(z)^m}{\mathcal V_{\alpha,m}(z)^m} = U (z)^m = \int_{z}^\infty Q(\tfrac{z}{\rho}) U(\rho) d \rho \to \frac{M}{2}\frac{\Gamma(b+1)}{\Gamma(b+1-\alpha)}.
\end{equation*}
\qed

%
\subsection{Kernel estimates}

We focus on the representation using \eqref{eq:Q}.
We begin this section with some notation
For $f, g: (a,b) \to [0,+\infty)$, introduce the notation
\[
  f \asympint{a}{b} g,
\]
if there exists $c(a,b)$, $C(a,b) > 0$ such that
\begin{equation*}
  c(a,b) g(z) \le f(z) \le C(a,b) g(z) \qquad \forall z \in (a,b).
\end{equation*}
We will set $\asymp$ to mean $\asympint{0}{+\infty}$.

\bigskip

For $\eta \sim 1^-$ (specifically for $\eta \in [1/2, 1)$), we have
\begin{equation*}
  \sigma^{1 + \frac1 b - d} \asympint{\frac 1 2}{1} 1 \qquad \text{and} \qquad 1-\sigma^{\frac 1 b} \asympint{\frac 1 2}{1} 1-\sigma.
\end{equation*}
Furthermore, because $\sigma \in (\eta, 1)$, as $\eta \to 1^-$ we also have $\sigma \to 1^-$, which implies $\eta\sigma^{-1} \to 1$.
Hence, we can estimate $Q(\eta)$ by focusing on the remaining singular term. Using the change of variable $(1-\eta)\xi = 1-\sigma$ (which gives $d\sigma = -(1-\eta)d\xi$), we recover:
\begin{equation*}
  \begin{aligned}
    Q(\eta) & \asympint{\frac{1}{2}}{1} \int_\eta^1 (1-\sigma)^{-\alpha} \, d\sigma
    \asympint{\frac{1}{2}}{1} \int_0^1 \big((1-\eta)\xi\big)^{-\alpha} (1-\eta) \, d\xi
    \asympint{\frac{1}{2}}{1} (1-\eta)^{1-\alpha} \int_0^1 \xi^{-\alpha} \, d\xi.
  \end{aligned}
\end{equation*}
We deduce that
\begin{equation}
  \label{eq:Q away from 0}
  Q(\eta) \asympint{\frac{1}{2}}{1} (1 - \eta)^{1-\alpha}
\end{equation}
We now estimate the behavior of $Q(0^+)$. When $0 < \eta < \frac 1 2$ then we write
\begin{align*}
  Q(\eta) & = \frac{1}{\Gamma(1-\alpha)}\left( \int_\eta^{\frac 1 2} (1-\sigma^{\frac 1 b})^{-\alpha} \sigma^{1-d} d\sigma +  \int_{\frac 1 2}^1 (1-\sigma^{\frac 1 b})^{-\alpha} \sigma^{1-d} d\sigma \right)
  \asympint{0}{\frac 1 2} \int_\eta^1 \sigma^{1-d} d \sigma + 1.
\end{align*}
Hence, we deduce that
\begin{equation}
  \label{eq:Q close to 0}
  Q(\eta) \asympint{0}{\frac 1 2}
  \begin{dcases}
    1          & \text{if } d = 1,   \\
    -\log \eta & \text{if } d = 2,   \\
    \eta^{2-d} & \text{if } d \ge 3.
  \end{dcases}
  \qquad
  .
\end{equation}
Hence, if $d \ge 2$ we have $K(0^+, \rho) = +\infty$ for each $\rho > 0$. This means that any non-trivial solution satisfies $U(0^+) = +\infty$.

%
\subsection{Limits of the kernel as $\alpha \to 1^-$}
Notice that
\begin{equation*}
  F_\alpha (z) = \int_0^1 \sigma^{-\alpha} U(\sigma^{-b} z) d \mu_\alpha(1-\sigma), \qquad d\mu_\alpha (z) = \frac{t^{-\alpha}}{\Gamma(1-\alpha)} dz.
\end{equation*}
It is easy to see that $\mu_\alpha \to \delta_0$ and hence
\begin{equation*}
  F_\alpha (z) \to U(z) \qquad \text{ as } \alpha \to 1^-.
\end{equation*}
Going back to \eqref{eq:implicit equation last formula with F} we can expect that
\begin{equation}
  \label{eq:K limit as alpha to 1}
  K_{\alpha,m}(z,\rho) \to K_{1,m} (\rho) \coloneqq b \rho \text{ as } \alpha \to 1 \text{ for each } 0 < z \le \rho.
\end{equation}
We prove this rigorously.
\begin{lemma}
  \label{eq:Q limit alpha -> 1}
  Assume $Q(\eta)$ vanishes for $\eta \ge 1$, and for $\eta \in (0,1)$ its derivative is given according to \eqref{eq:Q'}.
  %
  %
  %
  %
  Then, we have that
  %
  \[
    \lim_{\alpha \to 1^-} Q_{\alpha,m}(\eta) = Q_{1,m} \coloneqq b_{1,m}  \qquad \text{ for all } \eta \in (0,1).
  \]
\end{lemma}

\begin{proof}
  For this proof we will denote $b_\alpha = b_{\alpha,m}$. Using the change of variables $s = 1 - \sigma^{\frac{1}{b_\alpha}}$, which implies $\sigma = (1-s)^{b_\alpha}$ and $d\sigma = -b_\alpha(1-s)^{b_\alpha-1} ds$, the integration limits change from $[\eta, 1]$ to $[1 - \eta^{1/b_\alpha}, 0]$. We then have
  \begin{align*}
    Q_{\alpha,m}(\eta) & = \int_0^{1-\eta^{1/b_\alpha}} \frac{s^{-\alpha}}{\Gamma(1-\alpha)} (1-s)^{b_\alpha(1-d)} b_\alpha(1-s)^{b_\alpha-1} ds \\
    & = \int_0^{1-\eta^{1/b_\alpha}} \frac{s^{-\alpha}}{\Gamma(1-\alpha)} \Psi_\alpha(s) ds 
    \qquad \text{ where } \Psi_\alpha(s) = b_\alpha (1-s)^{b_\alpha(2-d)-1}.
  \end{align*}
  We decompose the integral into two parts:
  \begin{equation*}
    \left| Q_{\alpha,m} (\eta) - \Psi_\alpha(0) \int_0^{1-\eta^{1/b_\alpha}} \frac{s^{-\alpha}}{\Gamma(1-\alpha)} ds \right|
    \le
    \int_0^{1-\eta^{1/b_\alpha}}
    \frac{s^{-\alpha}}{\Gamma(1-\alpha)} \left| \Psi_\alpha(s) - \Psi_\alpha(0) \right|
    ds.
  \end{equation*}
  Observing that $\Psi_\alpha(0) = b_{\alpha,m}$, the first integral evaluates exactly to:
  \begin{align*}
    \Psi_\alpha(0) \int_0^{1-\eta^{1/b_\alpha}} \frac{s^{-\alpha}}{\Gamma(1-\alpha)} ds
    &=
    b_{\alpha,m} \frac{(1-\eta^{1/b_\alpha})^{1-\alpha}}{(1-\alpha)\Gamma(1-\alpha)}
    =
    b_{\alpha,m} \frac{(1-\eta^{1/b_\alpha})^{1-\alpha}}{\Gamma(2-\alpha)}.
  \end{align*}
  As $\alpha \to 1^-$, we have $b_{\alpha,m} \to b_{1,m}$, the exponent $(1-\alpha) \to 0^+$, and $\Gamma(2-\alpha) \to \Gamma(1) = 1$. Since $\eta \in (0,1)$, the base $(1-\eta^{1/b_\alpha})$ is strictly positive, meaning $(1-\eta^{1/b_\alpha})^{1-\alpha} \to 1$. Thus, this leading term converges exactly to $b_{1,m}$.

  For the second term, we point out that $\Psi_\alpha$ is smooth near $s=0$. Because $\eta \in (0,1)$, the upper integration limit is uniformly bounded away from $1$: there exists an $\varepsilon > 0$ such that $1-\eta^{1/b_\alpha} \le 1 - \varepsilon$ for all $\alpha$ close to $1$. On the interval $[0, 1-\varepsilon]$, $\Psi_\alpha$ is Lipschitz continuous uniformly in $\alpha$. Thus, there exists a constant $C > 0$ independent of $\alpha$ such that $|\Psi_\alpha(s) - \Psi_\alpha(0)| \le C s$.
  Therefore, we can bound the remainder by:
  \begin{align*}
    \int_0^{1-\varepsilon} \frac{s^{-\alpha}}{\Gamma(1-\alpha)} \left| \Psi_\alpha(s) - \Psi_\alpha(0) \right| ds
    & \le
    \frac{C}{\Gamma(1-\alpha)} \int_0^{1-\varepsilon} s^{1-\alpha} ds
    \le \frac{C}{\Gamma(1-\alpha)(2-\alpha)}.
  \end{align*}
  As $\alpha \to 1^-$, we have $\Gamma(1-\alpha) \to \infty$. This causes the upper bound to vanish. Combining both parts, we conclude that $\lim_{\alpha \to 1^-} Q_{\alpha,m}(\eta) = b_{1,m}$.
\end{proof}

\begin{remark}
  Notice that the convergence is not uniform in $[0,1]$ since $Q_{\alpha,m}(1) = 0$ for all $\alpha < 1$.
  A similar argument shows that $-Q'_{\alpha,m} \to b_{1,m} \delta_1$ in $\mathcal D'(0,1)$.
\end{remark}

%

\section{Linear diffusion $m = 1$}
\label{sec:m=1}

\subsection{Proof of \Cref{thm:existence} for $m = 1$}
\label{sec:existence proof m=1}
\paragraph{Solution by Fourier transform.}
The existence of a profile is well-known, see \cite{MR1061448, MR2047909,MR3631303}. In this case, we can take the Fourier transform of \eqref{eq:main} to deduce that
\begin{equation*}
  \partial_t^\alpha \widehat u(t, \xi) = -|\xi|^2 \widehat u(t,\xi).
\end{equation*}
Hence, the solution with general initial datum $u(0,x) = u_0(x)$ can be recovered by separation of variables $\widehat u(t,\xi) = \widehat Z(t,\xi) \widehat u_0(\xi)$
where for each $\xi$ fixed $\widehat Z(t,\xi)$ is the solution to
\[
  \begin{dcases}
    \partial_t^\alpha \widehat Z(t,\xi) = -|\xi|^2 \widehat Z(t,\xi), & \text{if } t > 0 \\
    \widehat Z(0,\xi)=1.
  \end{dcases}
\]
We have precisely the following
\begin{equation*}
  \widehat Z(t, \xi) = E_\alpha(-|\xi|^2t^\alpha), \qquad E_\alpha(s) = \sum_{k=0}^\infty \frac{s^k}{\Gamma(\alpha k +1)}.
\end{equation*}
where $E_\alpha$ is the Mittag-Leffler function.
And this yields the following
\[
  u(t,x) = \int_{\Rd} Z(t,x-y) u_0(y) dy.
\]
The function $Z$ is the so-called \emph{fundamental solution}.

\paragraph{Reference profile.}
$Z$ is a classical solution for $x \ne 0$ and $t > 0$. We take $\mathcal U_{\alpha,1}(z) = Z(1,z e_1)$ and this is a solution to \eqref{eq:profile equation}.
This already has mass $1$. And in \eqref{eq:rescaling profile to correct mass} we get $A = M$ and
\begin{equation*}
  U_{\alpha,m,M} (z) = M \mathcal U_{\alpha,1}(z).
\end{equation*}
We can apply \Cref{lem:sss is weak solution} and \Cref{lem:limit of flux holds} to conclude the result.
\qed

\subsection{Uniqueness. Proof of \Cref{thm:m=1 uniqueness}}
Let $u_1, u_2$ be two solutions to \eqref{eq:very weak solution}.
The difference $w = u_1 - u_2$ satisfies
\begin{equation*}
  \int_{0}^T  \int_{\Rd} w(t,x) \dRL \varphi(T-t,x)dxdt
  = \int_0^T \int_{\Rd} w(t,x) \Delta\varphi(T-t,x)  dx dt
\end{equation*}
To prove uniqueness, we use as test function in the weak formulation the solution to the following adjoint problem
\begin{equation*}
  \dRL \varphi_\varepsilon(t,x) - \Delta \varphi_\varepsilon(t,x) = f_\varepsilon(t,x) , \qquad \varphi(0,x) = 0,
\end{equation*}
for a function $f_\varepsilon$ that we construct in the following. Notice that since starting from $0$ initial datum this problem coincides with the Caputo one.

We construct $f_\varepsilon$ as follows. Fix a non-negative mollifier $\zeta \in C^\infty_c (\mathbb R^{d+1})$ with $\int_{\mathbb R^{d+1}} \zeta(t,x) dt dx = 1$. Define $\zeta_\varepsilon(t,x) = \varepsilon^{-(d+1)} \zeta( t / \varepsilon , x / \varepsilon)$.
For $0 < R_1 < R_2$ and $\varepsilon > 0$ we construct
\begin{equation*}
  f_\varepsilon(t,x) = t^2 \int_{0}^T \int_{B_{R_2} \setminus B_{R_1}} \zeta_\varepsilon(t-s,y-x) g(T-s,y) dy dx, \qquad g(t,x) =  \operatorname{sign} w(t,x)
\end{equation*}
Due to the slow growth in $t$ of $f$ at $t = 0$ we have that $\varphi_\varepsilon \in X \coloneqq C^1((0,T) \times \mathbb R^d) \cap C((0,T); C^2(\mathbb R^d))$. Since test functions in \eqref{eq:very weak solution} are required to be compactly supported, we can approximate $\varphi_\varepsilon$ in $X$  by a sequence of compactly supported functions. Hence $\partial_t^\alpha \varphi_\varepsilon$ and $\Delta \varphi_\varepsilon$ are approximated uniformly in compacts with uniform bounds $L^\infty((0,T) \times \mathbb R^d))$. Since $w \in L^1((0,T) \times \mathbb R^d)$ the weak formulation holds of $\varphi_\varepsilon$ using the dominated convergence theorem. Due to the weak formulation, we have that
\begin{equation*}
\int_0^T \int_{\mathbb R^d} w(t,x) f_\varepsilon(T-t,x) dx dt = 0.
\end{equation*}
Then letting $\varepsilon \to 0$
we deduce that
\begin{equation*}
\int_0^T \int_{B_{R_2} \setminus B_{R_1}} (T-t)^2 |w(t,x)| = 0.
\end{equation*}
Since this works for $R_1, R_2$ general, $w = 0$ and the proof is complete.

Given two profiles $U_1, U_2$ of mass $M$, we can use \Cref{lem:sss is weak solution} to deduce that the associated solutions $u_1, u_2$ are weak solutions. Therefore, we deduce $u_1 = u_2$ and therefore $U_1 = U_2$.
\qed

%

\section{Slow diffusion $m>1$}
In this setting, we will show that the profiles are compactly supported and therefore all the computations for the derivation of \eqref{eq:profile equation} are justified.
\subsection{Uniqueness of solutions to the profile equation}
Let
\begin{equation*}
\fixedpoint U (r) = \left( \int_{r}^\infty K(r, \rho) U(\rho) d\rho  \right)^{1/m}.
\end{equation*}
We have the following fundamental result.
\begin{theorem}[Forward uniqueness]
\label{thm:profile forward uniqueness}
Assume that $U_1 \in C((0,\infty))$ with $\text{supp}(U_1) = [0,R]$ and $U_2 \in C((0,\infty))$ are solutions of $U_i = \fixedpoint U_i$ in $(z_0, +\infty)$. If $U_1(z_0) = U_2(z_0)$ then $U_1 = U_2$ in $[z_0,+\infty)$.
\end{theorem}

%
%
%
%
%
%
%
%
%
%

We start with a lemma that strongly uses the homogeneity of the operator.
\begin{lemma}
\label{lem:comparison by homogeneity}
Assume $U_1, U_2 \in C([z_0,\infty))$ are non-negative, such that $U_i = \fixedpoint U_i$ in $[z_0,+\infty)$, and the set
\[
  \Lambda = \{ \lambda > 0 : \lambda U_1 \le U_2 \text{ in } [z_0,R] \}
\]
is non-empty, then $U_1 \le U_2$.
\end{lemma}
\begin{proof}
Consider $\lambda_* = \sup \Lambda > 0$. Since $K \ge 0$ we have that
\begin{equation*}
  U_2 = \fixedpoint U_2 \ge \fixedpoint (\lambda_* U_1) = \lambda_*^{\frac 1 {m}} \fixedpoint U_1 = \lambda_*^{\frac 1 {m}} U_1
\end{equation*}
Therefore, $\lambda_*^{\frac 1 {m}} \in \Lambda$, we deduce $\lambda_* \ge \lambda_*^{\frac 1 {m}}$, and hence $\lambda_* \in \{0\} \cup [1,+\infty)$. We conclude that $\lambda_* \ge 1$ and so $U_1 \le U_2$.
\end{proof}

%
%
%
%
%
%
%
We now proceed to the proof of the main result.
\begin{proof}[Proof of \Cref{thm:profile forward uniqueness}]
First, let us prove that $U_1$ and $U_2$ have the same support. Since $K$ is positive then the supports are of the form $[0,R_i]$ with $R_1 \in (z_0,\infty)$ and $R_2 \in [z_0,+\infty]$, and $U_i > 0$ in $(z_0,R_i)$.

Assume, towards a contradiction, that $R_1 < R_2$. Then, we have $U_2/U_1 \ge c > 0$ is bounded below in $\text{supp}(U_1)$.
Since outside $\text{supp}(U_1)$ the claim $\lambda U_1 \le U_2$ is trivial, then $c \in \Lambda$. Therefore, by the previous lemma $U_1 \le U_2$. Assume that $U_1 < U_2$ in some set $\omega \subset [z_0, R]$ of positive measure. Then, since $K(z_0, \rho) > 0$ for $\rho \in (z_0,+\infty)$,
$$
U_1(z_0) = \fixedpoint U_1 (z_0) < \fixedpoint U_2 (z_0) = U_2 (z_0).
$$
But $U_1(z_0) = U_2(z_0)$. We conclude that $U_1 = U_2$, and hence $R_1 = R_2$. This is a contradiction.

If $R_2 < R_1$ then $R_2$ is finite, and we can switch the roles in the previous argument. Hence, we deduce that $R_1 = R_2$.

Now let us consider $ U_2^{(\varepsilon)}$ a re-scaling of support $[0,(1+\varepsilon)R]$ using \eqref{eq:scaling solutions}.
%
%
%
Since $U_2^{(\varepsilon)}$ is non-increasing, and it is positive at $z = R$, then it is strictly positive in $[z_0,R]$. Since $U_1$ is bounded in $[z_0,R]$, then the set $\Lambda$ in \Cref{lem:comparison by homogeneity} is non-empty. And $U_1 \le U_2^{(\varepsilon)}$. Letting $\varepsilon \to 0$ we deduce $U_1 \le U_2$. By switching indices and repeating the argument $U_1 = U_2$.
\end{proof}

\begin{corollary}
\label{thm:comparison principle}
Let $U_1, U_2 \in C((0,\infty))$ be non-negative solutions with a bounded support of $U_i = \fixedpoint U_i$ in $(0,+\infty)$. If $U_1(z_0) < U_2(z_0)$ for some $z_0 > 0$ then for each $z\in (0, \infty)$ either $U_1(z) < U_2(z)$ or $U_1(z) = U_2(z) = 0$.
\end{corollary}

\begin{proof}
If the set $\{ z > 0 : U_1(z) = U_2(z) \}$ is empty, the proof is complete since $U_1$ and $U_2$ are continuous. Otherwise, let $z_1 = \inf\{ z > 0 : U_1(z) = U_2(z) \}$. By continuity $U_1(z_1) = U_2(z_1)$. By \Cref{thm:profile forward uniqueness}, $U_1 = U_2$ in $[z_1,+\infty)$.

If $U_2(z_1) = 0$, since $U_i$ is non-increasing we have that $U_1 = U_2 = 0$ in $(z_1, +\infty)$.
Therefore $z_0 < z_1$ and, by continuity,
we conclude that $U_2 > U_1$ in $(0,z_1)$, so the result is true.

Assume, towards a contradiction that $U_1(z_1) = U_2(z_1) > 0$.
First, we point out that $U_1/U_2 = 1$ in $[z_1,+\infty)$.
If $z_0 > z_1$ the proof we have a contradiction and the proof is complete.
If $z_0 < z_1$, since $U_1$ is non-increasing, it is bounded below by a positive constant in $[z_0,z_1]$.
And $U_2$ is bounded above and below in $[z_0,z_1]$.
Hence $U_1/U_2$ is bounded below $[z_0, +\infty)$.
We can therefore apply \Cref{lem:comparison by homogeneity} to deduce $U_2 \le U_1$ in $[z_0, +\infty)$.
This is a contradiction.
\end{proof}

\subsection{A reference shape for the solution of support $[0,1]$}
We define the reference shape functions
\begin{equation}
\mathcal V_{\alpha,m} (z) =
\begin{dcases}
  \text{given by } \eqref{eq:profile at z=0} & \text{if } z \in (0,1/2) \\
  (1-z)_+^{\frac{2-\alpha}{m-1}}             & \text{if } z \ge 1/2.
\end{dcases}
\end{equation}
Notice that $\mathcal V_{\alpha,m}$ is non-increasing. We will now show that for suitable $\overline C$ we have that $\overline{\mathcal U}(z) = \overline C \mathcal V_{\alpha,m}(z)$ is a supersolution.

\begin{lemma}
Let $\alpha \in (0,1)$, $m > 1$, and $d \ge 1$. We have $\fixedpoint \mathcal V_{\alpha,m} \asymp \mathcal V_{\alpha,m}$.
\end{lemma}

\begin{proof}
We separate the proof into different ranges of $z$. We will show only the proof of the upper bounds, and the lower bounds are proved similarly.

\subparagraph{Range $z \in (\frac 1 4,1)$ and $d \ge 1$.}
We have $K(z,\rho) = \rho Q(z/\rho)$ and hence using \eqref{eq:Q away from 0} we have that using the change of variable $(1-z) \sigma = 1-\rho$
\begin{align*}
  \int_z^1 K(z,\rho) \mathcal V_{\alpha,m} (\rho) d \rho
  & =   \int_z^1 \rho Q(z/\rho) (1-\rho)_+^{\frac{2-\alpha}{m-1}} d \rho
  \\
  & \le  C(\alpha,m,d)   \int_z^1 \rho (1-z/\rho)^{1-\alpha} (1-\rho)_+^{\frac{2-\alpha}{m-1}} d \rho
  \\
  %
  %
  %
  %
  & \le  C(\alpha,m,d) (1-z)^{1-\alpha + \frac{2-\alpha}{m-1} + 1} \int_0^1 (1 - \sigma)^{1-\alpha} \sigma^{\frac{2-\alpha}{m-1}} d \sigma
  %
  %
  \\
  & \le {C(\alpha,m,d) } \mathcal V_{\alpha,m}(z)^m.
\end{align*}
And the lower bound follows similarly bounding $\rho \ge z \ge \frac 1 4$.

\subparagraph{Range $z \in (0,1/4)$ when $d = 1$.}
We have that
\begin{align*}
  \int_z^1 K(z,\rho) \mathcal V (\rho) d \rho
  %
  %
  & \le C \left(  \int_z^{2z}  Q(z/ \rho) \rho d \rho + \int_{2z}^{1/2}   Q(z/\rho) \rho d \rho + \int_{1/2}^{1}  Q(z/\rho) (1-\rho)^{\frac{2-\alpha}{m-1}} d \rho
  \right)
  \\
  & \le C\left( \int_z^{2z}  (1 - z/\rho)^{1-\alpha}  \rho d \rho + \int_{2z}^{1/2}\rho d \rho+ \int_{1/2}^{1} (1-\rho)^{\frac{2-\alpha}{m-1}} d \rho
  \right)
  \\
  & \le C \left(
    z^2
  +  1 \right)
  %
  %
  \le C \mathcal V(z)^m.
\end{align*}
The lower bound follows equivalently.

\subparagraph{Range $z \in (0,1/4)$ when $d \ge 3$.}
We split the computation into three parts
\begin{align*}
  \int_z^1 K(z,\rho) \mathcal V_{\alpha,m} (\rho) d \rho
  & \le C\left( \int_z^{2z}  K(z, \rho) \rho^{-\gamma} d \rho
    + \int_{2z}^{1/2}   K(z, \rho) \rho^{-\gamma} d \rho
  + \int_{1/2}^{1}  K(z, \rho) (1-\rho)^{\frac{2-\alpha}{m-1}} d \rho \right).
\end{align*}
Now we deal with each term separately
\begin{align*}
  \int_z^{2z}  K(z, \rho) \rho^{-\gamma} d \rho
  & \le C \int_z^{2z}  Q(z/ \rho) \rho^{1-\gamma} d \rho
  \le C\int_z^{2z}  (1 - z/\rho)^{1-\alpha}  \rho^{1-\gamma} d \rho
  %
  %
  %
  %
  \le C z^{2-\gamma}.
\end{align*}
We also have that
\begin{align*}
  \int_{2z}^{1/2}   K(z, \rho) \rho^{-\gamma} d \rho
  & \le C \int_{2z}^{1/2}   Q(z/\rho) \rho^{1-\gamma} d \rho
  \le C \int_{2z}^{1/2}(z/\rho)^{2-d}\rho^{1-\gamma} d \rho
  %
  %
  %
  \le C z^{2-d} .
\end{align*}
Lastly
\begin{align*}
  \int_{1/2}^{1}  K(z, \rho) (1-\rho)^{\frac{2-\alpha}{m-1}} d \rho
  & \le C \int_{1/2}^{1}  Q(z/\rho) (1-\rho)^{\frac{2-\alpha}{m-1}} d \rho
  %
  %
  \le C \int_{1/2}^{1} (z/\rho)^{2-d} (1-\rho)^{\frac{2-\alpha}{m-1}} d \rho
  %
  %
  %
  %
  \le Cz^{2-d}.
\end{align*}
Because $m > 1$, we have $\gamma < d-2$, which implies that the term $z^{2-d}$ strictly dominates the term $z^{2-\gamma}$ as $z \to 0^+$. Given that $\gamma m = d-2$, this yields the correct target scaling. We conclude that $(\fixedpoint\mathcal{V}_d(z))^m \le C z^{2-d} \le C \mathcal{V}_d(z)^m$. The lower bound follows analogously.

\subparagraph{Range $z \in (0,1/4)$ when $d = 2$.}
We proceed in a similar way, computing each term
\begin{align*}
  \int_z^1 & K(z,\rho) \mathcal V_{\alpha,m} (\rho) d \rho
  \\
  & \le C \left( \int_z^{2z}  K(z, \rho) (-\log \rho)^{\frac 1 m} d \rho + \int_{2z}^{1/2}   K(z, \rho) (-\log \rho) ^{\frac 1 m}d \rho + \int_{1/2}^{1}  K(z, \rho) (1-\rho)^{\frac{2-\alpha}{m-1}} d \rho \right).
\end{align*}
We now estimate the terms separately. First, we deal with the influence of $\rho$ closest to $z$
\begin{align*}
  \int_z^{2z}  K(z, \rho) (-\log \rho)^{\frac 1 m} d \rho
  & \le C \int_z^{2z}  Q(z/ \rho) \rho(-\log \rho)^{\frac 1 m} d \rho
  \\
  & \le C  \int_z^{2z}  (1 - z/\rho)^{1-\alpha}  \rho(-\log \rho)^{\frac 1 m} d \rho
  \\
  & \le C z^2 \int_1^{2}  (1 - \sigma^{-1})^{1-\alpha}  \sigma (-\log z - \log \sigma)^{\frac 1 m} d \sigma
  \\
  & \le  C z^2 (-\log z)^{\frac 1 m} \int_1^{2}  (1 - \sigma^{-1})^{1-\alpha}  \sigma  d \sigma.
\end{align*}
Now we deal with the influence of intermediate values of $\rho$
\begin{align*}
  \int_{2z}^{1/2}  K(z, \rho) (-\log \rho)^{\frac 1 m} d \rho
  & \le \int_{2z}^{1/2}   Q(z/\rho) \rho(-\log \rho)^{\frac 1 m} d \rho
  %
  %
  \le C \int_{2z}^{1/2}(-\log z/\rho)\rho(-\log \rho)^{\frac 1 m} d \rho
  \\
  & \le C\left( (-\log z) \int_{z}^{\frac 1 2} \rho (-\log \rho)^{\frac 1 m} d\rho -   \int_{z}^{\frac 1 2} \rho (-\log \rho)^{\frac 1 m + 1}d\rho \right)
  \\
  & \le C (-\log z).
\end{align*}
Lastly, we control the terms where $\rho$ is very far from $z$
\begin{align*}
  \int_{\frac 1 2}^{1}  K(z, \rho) (-\log \rho)^{\frac 1 m} d \rho
  & \le \int_{1/2}^{1}  Q(z/\rho) (1-\rho)^{\frac{2-\alpha}{m-1}} d \rho
  %
  %
  \le C \int_{1/2}^{1} (-\log z/\rho) (1-\rho)^{\frac{2-\alpha}{m-1}} d \rho
  \\
  & 
  \le  C (-\log (z/2)) \int_{1/2}^{1} (1-\rho)^{\frac{2-\alpha}{m-1}} d \rho
  %
  %
  \le C (-\log z).
\end{align*}
Since $m > 1$, we conclude that $(\fixedpoint\mathcal V (z))^{m} \le C (-\log z) \le C \mathcal V(z)^m$. 
\end{proof}

%
%
%
%
%
%
%
%
%
%
%
%
%
%
%
%
%
%
%
%
%
%
%
%
%
%
%
%
%
%
%
%
%
\subsection{Proof of \Cref{thm:existence} for $m > 1$ and \Cref{thm:slow-diffusion}}
\label{sec:existence proof m>1}
\paragraph{Construction of the canonical profile of support $[0,1]$.}
Taking into account this last equivalence, we can construct the supersolutions
\begin{equation}
\label{eq:sub and supersolutions}
\begin{aligned}
  \overline{\mathcal U}(z)  & = \overline c \mathcal V_{\alpha,m}(z) \qquad \text{where } \overline c \ge \left(\sup_{z \in (0,1)} \frac{\fixedpoint \mathcal V_{\alpha,m} (z)}{\mathcal V_{\alpha,m}(z)}\right)^{\frac m {m-1}},
  \\
  \underline{\mathcal U}(z) & = \underline c \mathcal V_{\alpha,m}(z) \qquad \text{where } \underline c \le \left(\inf_{z \in (0,1)} \frac{\fixedpoint \mathcal V_{\alpha,m} (z)}{\mathcal V_{\alpha,m}(z)}\right)^{\frac m {m-1}}.
\end{aligned}
\end{equation}
We use the method of sub- and supersolutions using \eqref{eq:sub and supersolutions}. Let $U^{(0)} = \underline{\mathcal U}$ and $U^{(n)} = \fixedpoint U^{(n-1)}$. In every iteration, since $K \ge 0$ we have that
\begin{equation*}
U^{(n)} \le \fixedpoint U^{(n)} \le \fixedpoint \overline{\mathcal U} \le \overline{\mathcal U}.
\end{equation*}
Therefore, $U^{(n)}$ is a point-wise non-decreasing sequence of continuous functions bounded above by $\overline{\mathcal U}$. Therefore, it has a point-wise limit $\mathcal U$. Since $\overline{\mathcal U}$ provides an integrable bound for the kernel, we can apply the dominated convergence theorem to deduce that $\mathcal U = \fixedpoint \mathcal U$. Finally, since $\mathcal U$ is the image of the integral operator $\fixedpoint$, we immediately deduce $\mathcal U^m \in C^1((0,\infty))$. Due to the sub- and supersolutions, the support of $\mathcal U$ is precisely $[0,1]$. Furthermore, due to the integrability of $\mathcal U$, we have that $z^{d-1} \mathcal U(z) \in L^1 (0,\infty)$.

\paragraph{Construction of a solution of given mass.}
Using \Cref{sec:scaling} we construct the solution of mass $M$ by \eqref{eq:rescaling profile to correct mass}. The uniqueness of $U_{\alpha,m,M}$ is due to \Cref{thm:comparison principle}.

\paragraph{The solution satisfies the weak formulation}
We can apply \Cref{lem:sss is weak solution}.

\paragraph{Sharp constant at the free boundary.}
\label{sec:sharp constant at free boundary}
Similarly to \eqref{eq:sub and supersolutions}, we can use sub- and supersolutions in $(z,1)$ to deduce that
\begin{equation}
\label{eq:slow sub and supersolutions near free boundary}
\left(\inf_{\xi \in (z,1)} \frac{\int_\xi^1 K(\xi, \rho) \mathcal V(\rho) d\rho}{\mathcal V(z)^m}\right)^{\frac 1 {m-1}} \mathcal V(z) \le U(z) \le \left(\sup_{\xi \in (z,1)} \frac{\int_\xi^1 K(\xi, \rho) \mathcal V(\rho) d\rho}{\mathcal V(z)^m}\right)^{\frac 1 {m-1}}  \mathcal V(z).
\end{equation}
To compute this limit directly, we set $\gamma = \frac{2-\alpha}{m-1}$, which implies $\gamma m = 2 - \alpha + \gamma$. We evaluate the integral using the change of variables $1-\rho = (1-z)\sigma$, which gives $d\rho = (1-z)d\sigma$:
\begin{equation}
\label{eq:slow sharp constant at free boundary deduction}
\begin{aligned}
  \lim_{z \to 1^-} \frac{\int_z^1 K(z, \rho) \mathcal V(\rho) d\rho} {\mathcal V(z)^m}
  &= \frac{b^\alpha}{\Gamma(2-\alpha)} \lim_{z \to 1^-} \frac{1}{(1-z)^{\gamma m}} \int_{z}^1 (1-\tfrac{z}{\rho})^{1-\alpha} (1-\rho)^\gamma d \rho
  \\
  &= \frac{b^\alpha}{\Gamma(2-\alpha)} \lim_{z \to 1^-} \frac{1}{(1-z)^{\gamma m}} \int_{0}^1 \Big(\frac{(1-z)(1-\sigma)}{\rho}\Big)^{1-\alpha} (1-z)^\gamma \sigma^\gamma (1-z) d \sigma
  \\
  &= \frac{b^\alpha}{\Gamma(2-\alpha)} \lim_{z \to 1^-} \frac{(1-z)^{2-\alpha+\gamma}}{(1-z)^{\gamma m}} \int_{0}^1 (1-\sigma)^{1-\alpha} \sigma^\gamma d \sigma
  \\
  &=\frac{b^\alpha}{\Gamma(2-\alpha)} \int_{0}^1 (1-\sigma)^{1-\alpha} \sigma^\gamma d \sigma = \frac{b^\alpha}{\Gamma(2-\alpha)} B(2-\alpha, 1 + \gamma) \\
  &= \frac{b^\alpha \Gamma(1 + \gamma)}{\Gamma(3-\alpha+\gamma)},
\end{aligned}
\end{equation}
since $\rho \to 1$ as $z \to 1^-$, and the powers of $(1-z)$ cancel because $\gamma m = 2-\alpha+\gamma$. Evaluating the final Beta function integral, we recover the exact value for the sharp constant \eqref{eq:slow sharp constant at free boundary}.

\paragraph{Classical limit as $\alpha \to 1^-$.}
Due to \eqref{eq:Lieb upper bound} and \eqref{eq:equicontinuity} we have that $\{\mathcal U_{\alpha,m} : \alpha < 1 \}$ is pre-compact in $C_{loc}((0,1])$.
Hence, there is a subsequence $\mathcal U_{\alpha_k,m}$ that converges to $\widehat{\mathcal U}_m$. Next we observe from the supersolutions that $\mathcal U_{\alpha,m,d}$ are uniformly controlled as $\alpha \to 1^-$. So, the support of the limit is still $[0,1]$. Moving to the limit of the equation, we deduce that $\widehat{\mathcal U}_m$ is a solution of the equation at $\alpha = 1$. Lastly, we use uniqueness to show that it is $\mathcal U_{m}$.

%
%
%
%
%
%
%
%

%
\paragraph{Mesa limit.} Lastly, we discuss the limit $m \to \infty$.

%
%
%
%
%
%
%
%
%
%
%
%
%
%

\subparagraph{Mesa limit of $\mathcal U_{\alpha,m}$.}
For $z > 0$ we go back to the estimate near the free boundary \eqref{eq:slow sharp constant at free boundary}, and we adapt it slightly to
\begin{align*}
& \left(\inf_{\xi \in (z,1)} C(\xi)\right)^{\frac 1 {m-1}} (1-z)^{\frac{2-\alpha}{1-m}}
\le \mathcal U_{\alpha,m}(z)
\le \left(\sup_{\xi \in (z,1)} C(\xi)\right)^{\frac 1 {m-1}} (1-z)^{\frac{2-\alpha}{1-m}},
\end{align*}
where
\[
C(z) \coloneqq \frac{\int_z^1 K(z, \rho) (1-\rho)^{\frac{2-\alpha}{1-m}} d\rho}{(1-z)^{\frac{2-\alpha}{1-m}m}}.
\]
Using the explicit form of $Q_m'$ we can deduce similarly to \eqref{eq:slow sharp constant at free boundary deduction} that
\begin{equation}
\label{eq:slow upper and lower bounds near z=1}
(b^{-\alpha} c(d,\alpha,z))^{\frac{1}{m-1}} (1-z)_+^{\frac{2-\alpha}{m-1}} \le \mathcal U_{\alpha,m}(z) \le (b^{-\alpha} C(d,\alpha,z))^{\frac{1}{m-1}} (1-z)_+^{\frac{2-\alpha}{m-1}}
\end{equation}
Returning to \eqref{eq:slow upper and lower bounds near z=1}
Given that $b_{\alpha,m}^{\frac 1 {m-1}} = \left(\frac{\alpha}{2+d(m-1)}\right)^{\frac 1{m-1}} \to 1$ is as $m \to \infty$ we deduce that
\begin{equation*}
\lim_{m\to \infty} \mathcal U_{\alpha,m} (z) =
\begin{dcases}
  1 & \text{if } 0 < z < 1, \\
  0 & \text{if } z \ge 1.
\end{dcases}
\end{equation*}

\subparagraph{Mesa limit of $U_{\alpha,m,M}$.} Now we look at the construction of the profile of mass $M$ deduced in \eqref{eq:rescaling profile to correct mass} and let $L = A^{-\frac{m-1}{2}}$.
Due to the point-wise convergence and the existence of super-solutions we can use the dominated convergence theorem to deduce that
\begin{equation*}
L(\alpha,m,d,M) \to \frac{dM}{|\partial B_1|} = \frac{M}{|B_1|}.
\end{equation*}

\subparagraph{Mesa limit of $u_{\alpha,m,M}$.} Now we notice that $b_{\alpha,m} \to 0$, therefore we have that
\begin{equation*}
u_{\alpha,m,M}(t,x) = t^{-b_{\alpha,m} d} U_{\alpha,m,M}(t^{-b_{\alpha m}}|x|) \to \frac{M}{|B_1|} \chi_{[0,1]}(|x|).
\end{equation*}
This concludes the proof. \qed

\begin{remark}
We point out that the stability of \eqref{eq:slow sharp constant at free boundary} is consistent with the limit as $\alpha \to 1^-$.
%
%
%
%
%
We have that
\begin{equation*}
  \lim_{\alpha \to 1^-} \frac{\Gamma(1 + \tfrac{2-\alpha}{m-1} )}{\Gamma(3-\alpha+\tfrac{2-\alpha}{m-1})} = \frac{\Gamma(1 + \frac{1}{m-1})}{\Gamma(2 + \frac{1}{m-1})} = \frac{1}{1 + \frac{1}{m-1}} = \frac{(m-1)}{m},
\end{equation*}
and hence for $z \sim 1$
\begin{align*}
  \mathcal U_{\alpha,m}(z) &\sim \left( \frac{b^\alpha \Gamma(1 + \tfrac{2-\alpha}{m-1} )}{\Gamma(3-\alpha+\tfrac{2-\alpha}{m-1})} (1-z)_+ \right)^{\frac{2-\alpha}{m-1}}
  \sim
  \left( \frac{b_{\alpha,m}^\alpha \Gamma(1 + \tfrac{2-\alpha}{m-1} )}{\Gamma(3-\alpha+\tfrac{2-\alpha}{m-1})} \frac{(1-z^2)_+}{2} \right)^{\frac{2-\alpha}{m-1}}
  \\
  &\to
  \left( \frac{b_{1,m}(m-1)}{2m} (1-z^2)_+ \right)^{\frac{1}{m-1}} = \mathcal U_m(z).
\end{align*}
\end{remark}

%
%
%
%
%

%
%

%
%
%
%
%
%

%
%
%
%
%
%
%
%
%
%
%

%
%
%
%
%
%
%
%
%
%
%
%
%

%
%
%

%

%
%
%
%
%

%

%
%
%
%
%

%
%
%
%
%
%
%
%
%
%
%

%

%

%

%

%
%
%
%
%
%

%

%
%

%
%
%
%
%

%

%

%

\begin{comment}
\begin{theorem}
    \label{thm:profile well-posedness intermediate value}
    For every $y_0, z_0 > 0$ there exists a unique $U \in C((0,\infty))$ solution of \eqref{eq:profile equation} such that $U(z_0) = y_0$, it is compactly supported and has finite mass.
    Furthermore, for each $R > 0$ there is a unique continuous solution of support $[0,R]$, and it is given precisely by
    $$
        U_R(z) = R^{-\frac 2 {m-1}} U_1 (\tfrac{z}{R}).
    $$
    where $U_1$ is a reference profile of support $[0,1]$.
    \todo[inline]{
        \color{teal}Should $U_1$ be written in a different font as above?
        \color{blue}
        The $\mathcal V_{\alpha,m}$ is the reference profile, which is not a solution. Whereas $U$ is the solution. I am not picky with the notation. Perhaps we can remove the subindex for $\mathcal V_{\alpha,m}$.

        \color{teal}
        Also, the exponent in $R$ should be negative.
        \color{blue}
        Agreed!
        \color{teal}
        Now, I understand. I think we can leave the notation as it is.
        \normalcolor
    }
\end{theorem}
This result is a consequence of the following existence result, the scaling properties, and the forward uniqueness result.
\begin{lemma}
    \label{lem:existence of Barenblatt of support 0 1}
    There exists a $U \in C((0,\infty))$ solution to \eqref{eq:profile equation} such that $\text{supp}(U) = [0,1]$. Furthermore, for some $\underline c_1, \overline c_1 > 0$ we have that
    \[
        \underline c_1(z_0) (1 - z)_+^{\frac{2-\alpha}{m-1}} \le  U(z) \le \overline  c_1(z_0) (1 - z)_+^{\frac{2-\alpha}{m-1}} ,  \qquad \text{ for all } 0 < z_0 \le z < \infty.
    \]
\end{lemma}
\begin{proof}
    We use the method of sub- and supersolutions using \eqref{eq:sub and supersolutions}. Let $U_0 = \underline{\mathcal U}$ and $U_n = \fixedpoint U_{n-1}$. In every iteration, we have that
    \begin{equation*}
        U_n \le \fixedpoint U_n = U_{n+1} \le \fixedpoint \overline{\mathcal U} = \overline{\mathcal U}.
    \end{equation*}
    \color{teal}
    \todo[color=blue!20]{I agree with the text}
    Therefore, $U_n$ is a pointwise non-decreasing sequence of continuous functions bounded above by $\overline{U}$. It has a pointwise limit $U$. Since $\overline{U}$ provides an integrable bound for the kernel, we can apply the dominated convergence theorem to deduce that $U = \fixedpoint U$. Finally, since $U$ is the image of the integral operator $\fixedpoint$, we immediately deduce $U \in C((0,\infty))$.
    \normalcolor
\end{proof}
\end{comment}
%
%
%
%
%
%
%
%
%
%
%
%
%
%
%
%
%
%
%
%
%
%
%
%
%
%
%
%
%
%
%
%
%
%
%
%
%
%
%
%
%
%
%
%
%
%
%
%
%
%

\begin{comment}
\begin{proof}[Proof of \Cref{thm:profile well-posedness intermediate value}]
    Existence follows by the previous lemma and scaling.

\end{proof}
\end{comment}

%
%
%
%
%
%
%
%
%
%
%
%
%
%
%
%
%
%
%
%
%
%
%
%
%
%
%
%
%
%
%
%
%
%
%
%
%
%
%
%
%
%
%
%
%
%
%
%
%
%
%
%
%
%
%
%
%
%
%
%
%
%
%
%
%
%
%
%
%
%
%
%
%
%
%
%
%
%
%
%
%
%
%
%
%
%
%
%
%
%
%
%
%
%
%
%
%
%
%
%
%
%
%
%
%
%
%
%
%
%
%
%
%
%
%
%
%
%
%
%
%
%
%
%
%
%
%
%
%
%
%
%
%
%
%
%
%
%
%
%
%
%
%
%
%
%
%
%
%
%
%
%
%
%
%
%
%
%
%
%
%
%
%

%

\subsection{Numerical analysis for $m > 1$}

We now restrict ourselves to piece-wise constant functions. Then we look for the fixed point $U = \mathcal K_h U$ where
\begin{equation}
(\fixedpoint_h U)_i \coloneqq  \left( \sum_{j=i}^{+\infty} U_j K_{ij} \right)^{\frac 1 m}, \qquad \text{ where }
K_{ij} = \int_{z_{j}}^{z_{j+1}} K(z_i, \rho) d\rho.
\end{equation}
In dimension $d = 1$ we can use the closed formula for $Q$ given by \eqref{eq:Q from incomplete gammas}. This formula is stable as $\alpha \to 1^+$. For $d \ge 2$ we can integrate by parts
\begin{equation}
\begin{aligned}
  K_{ij} & = \int_{z_j}^{z_{j+1}} K(z_i, \rho) d\rho = \int_{z_j}^{z_{j+1}} \rho Q(z_i/\rho) d\rho = \color{teal}-\normalcolor\int_{z_j}^{z_{j+1}}  \int_{z_i/\rho}^1 \rho Q'(\eta)  d \eta d \rho
  \\
  & = \color{teal}-\normalcolor\int_{\frac{z_i}{z_{j+1}}}^{1} Q'(\eta) \frac{z_{j+1}^2 - \max\{ z_j, \frac{z_i}{\eta} \}^2} 2 d \eta
  \\
  & = \int_{\frac{z_i}{z_{j+1}}}^{1} \frac{\eta^{2-d-\frac{1}{b}}}{\Gamma(1-\alpha)} \frac{z_{j+1}^2 - \max\{ z_j, \frac{z_i}{\eta} \}^2} 2 (1-\eta^{1/b})^{-\alpha} \eta ^{1/b-1} d \eta,
  %
  %
\end{aligned}
\end{equation}
where we have used the change of variable $w = (1 - \eta^{1/b})^{1-\alpha}$. In the limit $\alpha \to 1^-$ we recover $K_{ij} \to b_1 \frac{z_{j+1}^2 - z_j^2}{2}$.
%
%
%
%
Our weights $K_{ij}$ our weights are not exactly computable. We have implemented this scheme in \texttt{julia} using the \texttt{Integrals.jl} package to compute these integrals, using the Gauss-Kronrod quadrature implemented in \texttt{QuadGK.jl}.
In $d=1$ our scheme is similar to \cite{Plociniczak2019}, where the effect of the quadrature formula for the computation of the weights.

\bigskip

We find approximate fixed points this problem by Picard iterations $U^{(n)} = \fixedpoint_h U^{(n-1)}$. We start iterating from a sub-solution we construct as follows.
We consider a mesh with $z_0 = 0$ and $z_I = 1$. As a initial datum we take $U^{(0)}_i = c \mathcal V_{\alpha,m}(z_i)$ with $c$ small enough
\begin{equation*}
c \le \left( \min_{i=1,\cdots,I-1} \frac{\sum_{j=i}^{I-1}  K_{ij} \mathcal V_{\alpha,m}(z_j) }{\mathcal V_{\alpha,m}(z_i)^m}\right)^{\frac 1 {m-1}}
\end{equation*}
so that $U^{(0)} \le \fixedpoint_h U^{(0)}$. This ensures that $U^{(n)} = \fixedpoint_h U^{(n-1)}$ is non-decreasing in $n$. Conversely, we can construct a supersolution by taking $\overline U_i = C \mathcal V_{\alpha,m}(z_i)$ with $C$ chosen large enough:
\begin{equation*}
C \ge \left( \max_{i=1,\cdots,I-1} \frac{\sum_{j=i}^{I-1}  K_{ij} \mathcal V_{\alpha,m}(z_j) }{\mathcal V_{\alpha,m}(z_i)^m}\right)^{\frac 1 {m-1}}.
\end{equation*}
This ensures that $U^{(n)} \le \overline U_i$.

\begin{remark}
From a theoretical point of view, if $\mathcal U$ is a subsolution to \eqref{eq:profile equation}, then we can take $U_i^{(0)} = \underline{\mathcal U}(z_i)$ with a non-increasing sub-solution we have that
\begin{align*}
  \underline{\mathcal U}(z_i)^m & \le  \int_{z_i}^{+\infty} K(z,\rho) \underline{\mathcal U}(\rho) d \rho   =  \sum_{j=i}^{+\infty}  \int_{z_j}^{z_{j+1}} K(z,\rho) \underline{\mathcal U}(\rho) d \rho
  \le \sum_{j=i}^{+\infty} \underline{\mathcal U}(z_j) K_{ij} .
\end{align*}
Therefore, the sequence $U^{(n+1)} = \fixedpoint_h U^{(n)}$ is monotone decreasing. This procedure does not allow us to construct supersolutions.
\end{remark}

\normalcolor

\begin{remark}
Convergence of this scheme as $h = \max_i |z_{i+1} - z_{i}| \to 0$ tends to $0$ can be accomplished proving Hölder estimates of the numerical solutions.
\end{remark}

%

%
%
%
%
%
%
%
%
%
%
%
%
%
%
%
%
%
%

\normalcolor

%
\section{Fast diffusion $m_c < m < 1$}

%
%
%
%
%
%
%
%
%
%

%
%
%
%
%
%
%
%
%
%
%
%
%
%
%
%
%
%

%
\subsection{Very singular solution}
Let us define the possibly infinite integral
\[
\mathsf Q(\gamma) \coloneqq \int_1^\infty Q(\sigma^{-1}) \sigma^{1-\gamma} d \sigma = \int_0^1 Q(\sigma) \sigma^{\gamma - 3} d \sigma.
\]
When $\gamma > \max\{2,d\}$, the boundary terms vanish, and we can integrate by parts to write:
\begin{equation}
\label{eq:mathsf Q}
\mathsf Q(\gamma) = \int_0^1 (-Q'(\sigma)) \frac{\sigma^{\gamma - 2}}{\gamma - 2} d \sigma = \frac{b}{\gamma - 2} \frac{B\Big( 1-\alpha, b(\gamma-d) \Big)}{\Gamma(1-\alpha)}= \frac{b}{\gamma - 2} \frac{\Gamma\Big( b(\gamma-d) \Big)}{\Gamma(1-\alpha + b(\gamma-d))}.
\end{equation}
Notice that if $m > m_c = \frac{d-2}{d}$, then $\gamma = \frac{2}{1-m} > \max\{2,d\}$.

\begin{lemma}
\label{lem:vss}
For $m \in (m_c,1)$ we have that \eqref{eq:vss} is a solution to \eqref{eq:profile equation}.
\end{lemma}
\begin{proof}
Consider a candidate of the form $U(z) = c z^{-\gamma}$. Then we have $z > 0$ (and putting $\rho=z/\sigma$)
\begin{align*}
  (\fixedpoint U(z))^m & = c \int_z^\infty \rho Q(z/\rho)\, {\rho ^{-\gamma}} \,d \rho       \\
  & =   c z^{2-\gamma} {\int_0^1 Q(\sigma) \sigma^{\gamma-3} d \sigma } \\
  & = c^{1-m} \int_0^1 Q(\sigma) \sigma^{\gamma-3} d \sigma
  \, (c z^{-\frac{2 - \gamma}m})^m.
\end{align*}
The correct power is therefore a $\gamma$ such that $-\gamma = \frac{2-\gamma}{m}$, that is, $\gamma = \frac{2}{1-m}$. Note that $\gamma>d$ exactly for $m> m_c=(d-2)/d$. We also have
%
\begin{equation*}
  { \int_0^1 Q(\sigma) \sigma^{\gamma-3} d \sigma \le C \left( \int_{1/2}^1  (1-\sigma)^{1-\alpha} \sigma^{\gamma-3} d \sigma + \int_0^{1/2} \sigma^{2-d} \sigma^{\gamma-3}\,d\sigma \right).}
\end{equation*}
Therefore, the integral is also finite if $m \in (1-\frac 2 d,1)$. We get an exact solution for the precise value
\[
  (c^*)^{m-1}= \int_0^1 Q(\sigma) \sigma^{\gamma-3} d \sigma= \mathsf Q(\tfrac{2}{1-m}).
\]
This concludes the proof.
\end{proof}

\subsection{Higher order expansion of the tails of solutions below \eqref{eq:vss}}

This is a decreasing and positive function.
By the dominated convergence theorem, we have $\mathsf Q(+\infty) = 0$. We have calculated $(c^*)^{1-m} = \mathsf Q(\frac{2}{1-m})$ and now we look for the first order expansion
\[
U(z) = c^* z^{-\frac 2 {1-m}} ( 1 + A z^{-\gamma} ).
\]
We directly compute that if $1 + Az^{-\gamma} \ge 0$ then as $z\rightarrow \infty$, we have\normalcolor
\begin{align*}
\frac{U(z)^m}{U^*(z)^m}                & = (1 + Az^{-\gamma})^m = 1 + m A z ^{-\gamma} + {{O(A^2z^{-2\gamma})}},
\\
\frac{(\fixedpoint U(z))^m}{U^* (z)^m} & = 1 + \frac{\mathsf Q(\frac{2}{1-m}+\gamma)}{\mathsf Q(\frac{2}{1-m})} A z^{-\gamma}.
\end{align*}
where $U^*$ is the VSS given by \eqref{eq:vss}.
%
Therefore, since $m \in (0,1)$ we can take $\gamma^*$ such that
\begin{equation}
\label{eq:gamma^star}
\mathsf Q\left(\tfrac{2}{1-m}+\gamma^*\right) = m \mathsf Q\left(\tfrac{2}{1-m}\right).
\end{equation}
And the constant $A$ is free.

\begin{remark}
We recall the explicit formula \eqref{eq:mathsf Q} in terms of functions allows to compute efficiently $\gamma^*$ using numerical methods.
\end{remark}

\paragraph{A sub-solution.}
%
%
%
%
%
%
%
%
%
%
%

%
To strictly estimate the error, recall that for $m \in (0,1)$, the map $s \mapsto s^m$ is strictly concave. Therefore, its tangent line at $s=1$ lies strictly above the graph, which gives the elementary inequality $(1-x)^m \le 1 - mx$ for all $x \in (0,1)$.
Setting $x = A z^{-\gamma^*}$ with $A > 0$, we immediately have:
\begin{equation*}
\frac{U(z)^m}{U^*(z)^m} = (1 - A z^{-\gamma^*})^m \le 1 - m A z^{-\gamma^*} = \frac{(\fixedpoint U(z))^m}{U^* (z)^m}.
\end{equation*}
Also, where $U(z) \le 0$, the positive part $(U(z))_+ = 0$ is trivially a subsolution. As a consequence, we have the following result.

\begin{lemma}
For $m \in (0,1)$ and $A\in \mathbb R$ we find that $\underline{\mathcal U}(z) = U^*(z) ( 1 - A z^{-\gamma^*})_+$ is a subsolution to \eqref{eq:profile equation} for all $z \in (0, +\infty)$.
\end{lemma}
%

\paragraph{A super-solution for the tails.}
To construct a supersolution, let us consider
\begin{equation}
\label{eq:m < 1 tail supersolution}
U(z) = U^*(z) \left( 1 - z^{-\gamma^*} + z^{-\frac{3}{2}\gamma^*} \right).
\end{equation}
Now we estimate
\begin{align*}
\frac{U(z)^m - (\fixedpoint U(z))^m}{U^*(z)^m}
& \ge \Lambda z^{-\frac{3}{2}\gamma^*}
%
%
+ R(z)  \quad \text{where } \Lambda =  m - \frac{\mathsf Q(\frac{2}{1-m}+\frac{3}{2}\gamma^*)}{\mathsf Q(\frac{2}{1-m})} \ge 0.
\end{align*}
where $R$ is the Taylor remainder.
\begin{lemma}
\label{lem:fast supersolution tails}
There exists $z_0(\alpha, m)$ such that
\begin{equation}
  \label{eq:fast supersolution tails}
  \overline{\mathcal U}(z) = U^*(z) (1 - z^{-\gamma^*} + z^{-\frac 3 2\gamma^*} ),
\end{equation}
is a super-solution for all $z \ge z_0(\alpha,m)$.
\end{lemma}
\begin{proof}
We point out that $\overline{\mathcal U}(z) / U^*(z) \in [1,2]$ for $z \ge z_0 \ge 1$ and therefore
the Taylor remainder term can be controlled by
\begin{align*}
  |R(z)| \le  \frac{1}{2}\left( z^{-\gamma^*} - z^{-\frac 3 2 \gamma^*}  \right)^2 \sup_{\xi \in [1,2]} |m(m-1) \xi^{m-2}| \le  \frac{|m (m-1)|}{2} z^{-2\gamma^*} .
\end{align*}
We choose $z_0$ such that $|R(z_0)| \le \Lambda z_0^{-\frac 3 2 \gamma^*}$.
\end{proof}
%

%
%
%
%
%
%
%
%
%
%
%
%
%
%
%
%
%
%
%
%
%
%
%
%
%
%
%
%
%
%
%
%
%
%
%
%
%
%
%
%
%
%
%
%
%
%
%
%
%
%
%
%
%

%

%
%
%
%
%
%

%
%
%
%
%
%
%
%

%
%
%
%
%
%
%
%
%
%
%
%
%
%

%

%
%
%
%
%

%

%
%

%

%

%

\subsection{Super-solutions close to $z = 0$}
%
%
\begin{lemma}
\label{lem:fast supersolution close 0}
Assume $d \ge 2$, $m > 1 - \frac{2}{d}$ and let
\begin{equation*}
  \gamma \in
  \begin{dcases}
    [\tfrac{d-2}{m}, d) & \text{if } d \ge 3, \\
    (0,2)               & \text{if } d= 2.
  \end{dcases}
\end{equation*}
Let us pick
\[
  c^{1-m} = \frac{(d-2)\Gamma(1-\alpha)}{1 + \int_z^1 (1-{\rho}^{-1/b})^{-\alpha} \rho^{-\gamma+d-1} d \rho }.
\]
Assume that there exists a solution $U$ such that
\begin{equation*}
  \int_1^\infty (1-\rho^{-\frac 1 b})^{-\alpha} U(\rho) \rho^{d-1} d \rho \le c.
\end{equation*}
Then, we have
\begin{equation*}
  \overline{ U}(z) =
  \begin{dcases}
    c z^{-\gamma} & \text{if } z < 1,
    \\
    U(z)          & \text{if } z \ge 1,
  \end{dcases}
\end{equation*}
is a super-solution and, furthermore, $U(z) \le \overline{ U}(z)$ for $z \in (0,+\infty)$.
\end{lemma}
%
\begin{proof}
For $z \ge 1$ we have that $U$ is exactly a solution. For $z \in (0,1)$
\begin{align*}
  \int_z^\infty & (-Q'(z/\rho) )\overline{ U} (\rho) d \rho
  \\
  & \le \frac{z^{1-d}}{\Gamma(1-\alpha)}\left( c\int_z^1 (1-\left(\tfrac{z}{\rho}\right)^{1/b})^{-\alpha} \rho^{-\gamma+d-1} d \rho +\int_1^\infty (1-\left(\tfrac{z}{\rho}\right)^{1/b})^{-\alpha} U(\rho) \rho^{d-1}d\rho \right)
  \\
  %
  %
  & \le \frac{z^{1-d}}{\Gamma(1-\alpha)} c  \left(\int_z^1 (1-\left(\tfrac{1}{\rho}\right)^{1/b})^{-\alpha} \rho^{-\gamma+d-1} d \rho + 1\right) =
  (d-2) c^m z^{1-d}
  \\
  & = -(\overline{ U}^m)'(z).
\end{align*}
Since $\underline{ U} = 0$ is a sub-solution, we have that $\widetilde U = \fixedpoint^n \overline{ U}$ is a solution to the problem. Given that $\overline{ U} (z) = U(z)$ for $z > 1$ we have that $\fixedpoint^n \overline{ U}(z) = U(z)$ for $z > 1$, and therefore
$\widetilde U(z) = U(z)$ for $z > 1$.
Consider
\[
  z_0 = \inf\{ z > 0 : U = \widetilde U \text{ in } (z,+\infty) \}.
\]
Assume that $z_0 > 0$. For $z < z_0$ we write using the intermediate value theorem
\begin{align*}
  |U(z) - \widetilde U(z)| & = \frac {\xi^{\frac 1 m - 1}}{m} \left| \int_z^{z_0} \rho Q(z/\rho) (U(\sigma) - \widetilde U(\sigma)) d \sigma \right|                                                          \\
  & \le \left(\sup_{\sigma \in [z,z_0]} |U(\sigma) - \widetilde U(\sigma)|\right) \frac{ U(z)^{\frac 1 m - 1} + \widetilde U(z)^{\frac 1 m - 1}} m \int_{z}^{z_0} \rho Q(z/\rho) d \rho.
\end{align*}
Now taking supremum again to some $z_1 < z_0$ we recover
\begin{align*}
  \sup_{z \in [z_1,z_0]} & |U(z) - \widetilde U(z)|
  \\
  & \le \sup_{z \in [z_1,z_0]} |U(z) - \widetilde U(z)|  \sup_{z \in [z_1,z_0]} \frac{U(z)^{\frac 1 m - 1} + \widetilde U(z)^{\frac 1 m - 1}}{m} \int_{z}^{z_0} \rho Q(z/\rho) d \rho.
\end{align*}
Notice that
\begin{equation*}
  \lim_{z_1 \to z_0^-} \sup_{z \in [z_1,z_0]} \frac{U(z)^{\frac 1 m - 1} + \widetilde U(z)^{\frac 1 m - 1}}{m} \int_{z}^{z_0} \rho Q(z/\rho) d \rho = 0.
\end{equation*}
For $z_1 < z_0$ but close enough to $z_0$ this is smaller than $0$, and hence we deduce that
\[
  \sup_{z \in [z_1,z_0]} |U(z) - \widetilde U(z)| = 0.
\]
But this is a contradiction to the definition of $z_0$.
\end{proof}

%

%
%
%
%
%
%
%
%
%
%
%
%
%
%
%
%
%
%
%
%
%
%
%
%
%
%
%
%
%
%
%
%
%
%
%
%

%
%
%
%
%
%
%
%
%
%
%

\subsection{Proof of \Cref{thm:existence} for $m \in (m_c,1)$ and \Cref{thm:fast-diffusion}}
\label{sec:existence proof m < 1}
\paragraph{Existence of a family of profiles.}
Using the subsolution $\underline{\mathcal U}(z) = c^*z^{-\frac{2}{1-m}}(1 - z^{-\gamma^*})_+$ and $\overline{\mathcal U} = U^* $ we deduce that
\begin{equation}
\label{eq:m<1 canonical solution}
\mathcal U = \lim_n \fixedpoint^n \underline{\mathcal U},
\end{equation}
is in $C((0,+\infty))$ and it is a solution to the problem.
Furthermore, using \eqref{eq:fast supersolution tails} we deduce that $\mathcal U \ne U^*$ and we have the estimate
\begin{equation*}\label{Lbeh.VSSfan}
U^*(z) ( 1 - z^{-\gamma^*} )_+
\le \mathcal U(z) \le
U^*(z) (1 - z^{-\gamma^*} + z^{-\frac 3 2 \gamma^*}), \qquad \text{for all } z \ge z_0 > 0.
\end{equation*}
Let us define
\[
U_R(z) \coloneqq R^{\frac{-2}{1-m}} \mathcal U\left(\frac{z}{R}\right).
\]
\begin{lemma}
\label{lem:fast rescaled limits in R}
For $m \in (0,1)$ and $z > 0$ we have that $U_R(z)$ is a decreasing sequence in $R$ and, furthermore, for each $z>0$,
\[
  U_R(z) \to
  \begin{dcases}
    U^*(z) & \text{as } R \to 0,      \\
    0      & \text{as } R \to \infty.
  \end{dcases}
\]
\end{lemma}
\begin{proof}
Notice that
\[
  U_R(z) = c^*z^{-\frac 2{1-m}}\Big(1 - R^{\gamma^*} z^{-\gamma^*} + o\left(R^{\gamma^*} z^{-\gamma^*}\right)\Big).
\]
Hence, for $z$ fixed, we have
\[
  U_R(z) \to U^*(z) \text{ as } R \to 0.
\]
Furthermore, if $\overline R > \underline R$ then we have that
$U_{\overline R}(z) < U_{\underline R}(z)$ for $z$ sufficiently large. Consider
\begin{equation*}
  z_1 = \inf\{Z > 0 : \text{for all } z \in (Z,+\infty) \text{ we have that } U_{\overline R}(z) < U_{\underline R}(z)\}.
\end{equation*}
Assume that $z_1 > 0$. Then $U_{\overline R}(z_1) = U_{\underline R}(z_1)$. However, we have
\[
  U_{\overline R}(z_1)^m  = \int_{z_1}^\infty K(z,\rho) U_{\overline R}(\rho) d\rho < \int_{z_1}^\infty K(z,\rho) U_{\underline R}(\rho) d\rho = U_{\underline R} (z_1)^m ,
\]
which is a contradiction. There $U_R(z)$ is decreasing in $R$ for all $z>0$. Now consider
\[
  U_\infty(z) = \lim_{R \to \infty} U_R (z).
\]
Due to the monotone convergence theorem $U_\infty(z) = T U_\infty(z)$ for each $z > 0$. Furthermore, we have
\begin{align*}
  U_\infty(z) & = \inf_{R > 0} U_R(z) = \inf_{R>0} R^{-\frac{2}{1-m}} \mathcal U(z/R) =  \inf_{R>0} U^*(z) \frac{\mathcal U(z/R)}{U^*(z/R)} = U^*(z) \inf_{R>0} \frac{\mathcal U(z/R)}{U^*(z/R)} \\
  & = U^*(z) \inf_{y > 0} \frac{\mathcal U(y)}{U^*(y)} = c_\infty U^*(z).
\end{align*}
Since $\mathcal U < U^*$ we have $c_\infty < 1$. Because $U_\infty$ satisfies $U_\infty = \fixedpoint U_\infty$, we have $c_\infty^m U^*(z)^m = c_\infty U^*(z)^m$, which requires $c_\infty^m = c_\infty$. Since $m \in (0,1)$ and $c_\infty < 1$, we deduce that $c_\infty = 0$.\normalcolor
\end{proof}

\paragraph{For $d=1$ solutions are bounded.}
%
%
%
%
Notice that since $U_R$ is decreasing we can use \eqref{eq:Q'} with $d=1$ and observe that $-Q'$ is increasing. Hence for $z \in (0,1)$ we have that
\begin{align*}
-(U_R^m)'(z) & = \int_z^\infty (-Q')\left(\tfrac{z}{\rho}\right) U_R(\rho) d \rho                                                              \\
& \le U_R(z) \int_z^1 (-Q')\left(\tfrac{z}{\rho}\right) d \rho + \int_1^\infty (-Q')\left(\tfrac{z}{\rho}\right) U_R(\rho) d \rho \\
& \le U_R(z) \int_z^1 (-Q')\left(\rho^{-1}\right) d \rho + \int_1^\infty (-Q')\left(\rho^{-1}\right) U_R(\rho) d \rho             \\
& \le U_R(z) A + B_R.
\end{align*}
Due to \Cref{lem:fast rescaled limits in R} and the dominated convergence theorem $B_R \to 0$ as $R \to \infty$. We also have $U_R(1) \to 0$. Let us define $W(s) = [U_R(1-s)]^m$. Then we have
\begin{equation*}
W'(s) \le A W(s)^{\frac 1 m} + B_R, \qquad W(0) = U_R(1)^m.
\end{equation*}
Integrating this differential inequality, we see that if $B_R$ and $U_R(1)$ are small enough, then $U(0)^m = W(1) < \infty$.
Since $U_R$ is bounded for $R$ large enough, then $\mathcal U$ is bounded. Due to tail estimates, $\mathcal U$ has finite mass.
%

\paragraph{For $d \ge 2$ solutions have finite mass.}
We can apply \Cref{lem:fast supersolution close 0} to $U_{R(c)}$ with $R(c)$ large enough. Since $U_{R(c)} \le c z^{-(d-2)/m}$ we conclude that $U_R$ has finite mass for all $R > 0$. And hence so that $\mathcal U$.

\paragraph{Conclusion for $\alpha$ fixed.}
Since we have proven that solutions have finite mass, we can apply \eqref{eq:limit of flux at z=0} to obtain their sharp behavior at $z = 0$. Through \eqref{eq:equation for dUm/dz} we deduce that $U^m \in C^1(0,+\infty)$ and hence we can apply \Cref{lem:sss is weak solution} to deduce that they are weak solutions.

\paragraph{Limit as $\alpha \to 1^-$.} We point out that $c^*_{\alpha,m}$ is continuous at $\alpha = 1$, ahd hence $U^*_{\alpha,m}$ is continuous at $m = 1$.
Due to \Cref{eq:Q limit alpha -> 1} we observe that if $\gamma > 2$
\begin{equation*}
\lim_{\alpha \to 1^-}\mathsf Q_{\alpha,m} (\gamma) = \frac{b}{\gamma - 2}.
\end{equation*}
Hence, we can pass the limit $\gamma^*_{\alpha,m}$ to the limit as $\alpha \to 1$. And thus we can pass to the limit $\underline{\mathcal U}_{\alpha,m}$ and $\overline{\mathcal U}_{\alpha,m}$. Following the choice in \Cref{lem:fast supersolution tails} we see that $z_0(\alpha,m)$ is bounded as $\alpha \to 1^-$.
Using \eqref{eq:Lieb upper bound} and \eqref{eq:equicontinuity} we have that $\mathcal U_{\alpha,m}$ can be precompact in $C_{loc}((0,+\infty))$. Any possible limit $\widehat {\mathcal U}$ satisfies $z \ge z_0(1,m)$ that
\begin{equation}
\label{eq:m < 1 alpha = 1 upper and lower bounds}
U^*_{1, m} (z) (1 - z^{\gamma^*_{1,m}})_+ \le \widehat {\mathcal U} (z) \le U^*_{1, m} (z) (1 - z^{\gamma^*_{1,m}} + z^{-\frac 3 2 \gamma^*_{1,m}})
\end{equation}
for $z$ large enough.
Furthermore, using super-solutions and the dominated convergence theorem, it is a solution to \eqref{eq:profile equation} for $\alpha = 1$. Since $\widehat{\mathcal U}$ is a solution of \eqref{eq:profile equation} with $K_{1,m}(z,\rho) = b \rho$ such that \eqref{eq:m < 1 alpha = 1 upper and lower bounds} we can take one derivative and prove that
\begin{equation*}
-(\widehat{\mathcal U}^m)'(z) = b z \widehat{\mathcal U}(z).
\end{equation*}
Integrating this equation, we deduce that
\begin{equation*}
\widehat{\mathcal U}(z) = c^*_{1,m}(C + z^2)^{-\frac{1}{1-m}}.
\end{equation*}
Due to the asymptotic expansion, we conclude that $C = 1$. Since every convergent sub-sequence has the same limit, the whole sequence $\mathcal U_{\alpha,m}$ converges, and this is its limit.
\qed

\section{Limit $m \to 1$. Proof of \Cref{thm:limit m->1}}
Like in the case $\alpha = 1$, we expect continuity with respect to $m$ and $U_{\alpha,m,M}$ with respect to $m$. Already in this case, the limits are tricky to compute. We have some preliminary results.
Due to \eqref{eq:Lieb upper bound}  and \eqref{eq:equicontinuity} we have that $\{ U_{\alpha,m,M}^m : m \ne 1 \}$ is pre-compact in $C_{loc}((0,+\infty))$.
Consider $\widehat U$ the limit of any convergent sub-sequence with $m_j \to 1$.
Also, by Fatou's lemma we have that
\begin{equation*}
\widehat M = |\partial B_1| \int_0^\infty \widehat U(z) z^{d-1} dz \le M .
\end{equation*}

%

%

%
%

%
%

%

%
%

%

%

%

%
%

%
%
%
%
%
%
%
%
%
%
%
%
%
%
%

%
Assume, in contradiction, that the tightness in $L^1(0,\infty, z^{d-1})$ does not hold. Then, there are some $c > 0$, and sequences $R_j \to \infty$ and a subsequence $m_j \to 1$ (which we do not relabel) such that $V_j = U_{\alpha,m_j,M}$
\begin{equation*}
\int_{R_j}^\infty V_j(\rho) \rho^{d-1} d \rho \ge c.
\end{equation*}
Then we have that
\begin{equation*}
z^{d-1}(-V_j^{m_j}(z))' \ge \int_z^\infty - \frac{z^{d-1}}{\rho^{d-1}}Q'(\tfrac{z}{\rho}) \rho^{d-1} \widehat{U}(\rho) d \rho \ge \frac{c}{\Gamma(1-\alpha)} \text{for } z < R_j.
\end{equation*}
Integrating once, we deduce that
\[
V_j^{m_j}(z) \ge \frac{c}{\Gamma(1-\alpha)} \int_{z}^{R_j} \rho^{1-d} d\rho \qquad \text{for all } z < R_j.
\]
\paragraph{Case $d=1$.} Then we have $V_j^{m_j}(z) \ge \frac{c}{\Gamma(1-\alpha)} (R_j - z) \to \infty$ for each $z$ fixed as $j \to \infty$. This contradicts the point-wise convergence.

\paragraph{Case $d=2$.} Then we have $V_j^{m_j}(z) \ge \frac{c}{\Gamma(1-\alpha)} \log \frac{R_j}{z} \to \infty$ for each $z$ fixed as $j \to \infty$. This contradicts the point-wise convergence.

\paragraph{Case $d \ge 3$.} Then we have that $V_j^{m_j}(z) \ge \frac{c}{\Gamma(1-\alpha)(d-2)}(z^{2-d} - R_j^{2-d})$.
This implies that the point-wise limit satisfies
\[
\widehat U (z) \ge \frac{c}{\Gamma(1-\alpha)(d-2)} z^{2-d}.
\]
Using Fatou's lemma, we deduce that
\[
\widehat M \ge |\partial B_1| \int_0^\infty \rho^{d-1} V(\rho) d\rho \ge \frac{c |\partial B_1|}{\Gamma(1-\alpha)(d-2)}\int_0^\infty \rho d \rho = +\infty.
\]
This is a contradiction. Hence, the family $z^{d-1} U_{\alpha,m_j,1}(z)$ is tight in $L^1(0,\infty)$.

Due to the tightness and the convergence over compactness, we deduce that as $j \to \infty$ we have that
\[
\int_0^\infty |U_{\alpha,m_j,M}(z) - \widehat U(z)| z^{d-1} dz \to 0.
\]
Due to the bounds on $Q_{\alpha,m}$ and the convergence $Q_{\alpha,m} \to Q_{\alpha, 1}$ as $m \to 1$ this is sufficient to prove that
\begin{equation*}
\widehat U(z) = \int_z^\infty K_{\alpha,1}(\rho) \widehat U(\rho) d \rho , \qquad |\partial B_1| \int_0^\infty \widehat U(\rho) \rho^{d-1} d \rho = M.
\end{equation*}
By the uniqueness result \Cref{thm:m=1 uniqueness} this implies that $\widehat U = U_{\alpha,1,M}.$
For every $m_j \to 1$ the sequence $U_{\alpha,m_j,M}$ is pre-compact in $C_{loc}((0,+\infty))$, and the limit of any convergent subsequence is $U_{\alpha,1,M}$, then the whole sequence converges.
\qed

\section{Extensions and open problems}

\paragraph{Uniqueness of weak solutions with Dirac initial data.}
We expect that the arguments in \cite{Pierre1982} for the uniqueness of the Porous-Medium Equation with measure initial data can be adapted to this setting. The idea of the proof is to study $v = E * u$ (where $E$ is either the Green function for $-\Delta$ if $d \ge 3$ or $-\Delta + \varepsilon I$ if $d =1,2$).

\paragraph{Asymptotic behavior for general initial data.}
We have constructed a family of self-similar profiles $U_{\alpha,m}$. 
An open problem remains regarding the stability of these profiles: are they long-time global attractors for the corresponding general Cauchy problem? In the classical case $\alpha=1$, this problem is very well understood \cite{vazquez2014barenblatt}. 
For $\alpha \in (0,1)$ and $m = 1$ this was shown in \cite{cortazarHeatEquationMemory2021}.
However, in the time-fractional model, there are no corresponding results in the literature. 

\paragraph{The very-singular fast-diffusion problem.}
The range $0<m<m_c$, $d\ge 3$, is usually called very singular fast diffusion. The tools used in this paper are not suitable for treating this range for the equation \eqref{eq:main} with the fractional time derivative. To our knowledge, the problem is completely open.

In the standard case $\alpha=1$ much is known. According to the theory developed, fundamental solutions with finite mass exist but do not have the type of self-similarity used in this paper, there is no mass conservation, and they extinguish in finite time: see \cite{chasseigne2002theory,Vazquez2006,Vazquez2006a, vazquez2017, bonforte2010sharp}. \normalcolor

\paragraph{Fully non-local problems.}
The Porous Medium in non-local pressure \cite{allen2016}.
Another interesting family of problems is $\partial_t^\alpha u = (-\Delta)^s u^m$. Let the linear case $m = 1$ already be well understood; see \cite{cortazarLargeTimeBehaviorFully2021}.

\paragraph{Numerical analysis for $m_c < m < 1$.}
In this paper, we include a discussion of the scheme for $m > 1$. This case is manageable, as we have compactly supported profiles. From a theoretical point of view, the scheme is also valid for $m_c < m < 1$. In this case, the solutions have power-type tails, and hence smart truncations are needed in order to obtain an implementable scheme.

\paragraph{The asymptotic fan for $0 < m < m_c$ above $U^*$.}
For $\alpha = 1$ it is know that there exists self-similar profiles of the form $U^*(z) ( 1 + z^{-\gamma} + o(z^{-\gamma}))$ that blow up at some $z = z_1$. By re-scaling this creates an ordered family of solutions above $U^*(z)$, each of which blow up at a different $z > 0$.

\paragraph{Uniqueness solutions to \eqref{eq:profile equation} when $m < 1$.}
In this manuscript we have only proven the uniqueness of solutions to \eqref{eq:profile equation} in the case $m > 1$, by taking advantage of the super-linear homogeneity and the compact support. The case $m_c < m < 1$ remains open.

\paragraph{Travelling waves.}  
This is a topic where much is known in the standard case and probably nothing in the fractional time derivative case.\normalcolor

\paragraph{Sign-changing self-similar profiles.} In principle, we cannot rule out the possibility that there exist sign-changing solutions to \eqref{eq:profile equation}. We do not find them particularly physically relevant and we will not discuss them further.

\section*{Acknowledgments}

The research of DGC is supported by grants RYC2022-037317-I and  PID2021-127105NB-I00 from Agencia Estatal de Investigación of the Spanish Government. \L{}P has been supported by the National Science Centre, Poland (NCN) under the grant Sonata Bis with a number NCN 2020/38/E/ST1/00153.
JLV was supported by CNS2024-154515 from AEI, the Agencia Estatal de Investigación of the Spanish Government.

\printbibliography

\appendix

%
%
%
%
%
%
%
%
%
%
%
%
%
%
%
%
%
%
%
%
%
%
%
%
%
%
%
%
%
%
%
%
%
%
%
%
%
%
%
%
%
%
%

%
%

%

%

%
%
%
%
%
%
%
%
%
%
%
%
%
%
%
%
%
%
%
%
%

%
%
%
%
%
%
%
%
%
%
%
%
%
%
%
%
%
%
%
%
%
%
%
%
%
%
%
%
%
%
%
%
%
%
%
%
%
%
%
%
%
%
%
%
%
%
%
%
%
%
%
%
%
%
%
%
%
%

%
%
%
%
%
%
%
%
%
%
%
%
%
%
%
%
%
%
%
%
%
%
%
%
%
%
%
%
%
%
%
%
%
%
%
%
%
%

%
%
%
%
%
%
%
%
%
%
%
%
%
%
%
%
%
%
%
%
%
%
%
%
%
%
%
%
%
%
%
%
%
%
%
%
%
%
%

%
%
%
%
%
%
%
%
%
%
%
%
%
%
%
%
%
%
%
%
%
%
%
%
%
%
%
%
%
%
%
%
%

%
%
%
%
%
%
%
%
%
%
%
%
%
%
%
%
%
%
%
%
%
%
%
%
%
%
%
%
%
%
%

%

%
%
%
%
%
%
%

%
%
%
%
%
%
%
%
%
%
%
%
%
%
%
%
%
%
%
%

%
%

%

%

%
%

%
%
%

%
%
%
%
%
%
%
%
%
%
%
%
%
%
%
%
%
%
%
%
%
%
%
%
%
%
%
%
%
%
%
%
%
%
%

%
%
%
%
%
%
%
%

%

%

%

%

%

%

%

%
%
%
%
%
%
%
%
%
%
%
%
%
%
%
%
%

%
%
%
%
%
%
%
%
%
%
%
%
%
%
%
%
%
%
%
%
%
%
%
%
%
%
%
%
%
%

%

%
%
%
%
%
%
%
%
%

%
%
%
%
%
%
%
%

%
%
%
%
%
%
%
%
%
%
%
%
%
%
%

%

%
%
%
%
%
%
%
%

%

%
%
%

%

%
%

%

%

%

%

%

%

%
%
%

%
%

%
%

%

%
%
%

%
%
%
%

%
%

%
%
%
%
%
%

%

%

\section*{Addresses}

\noindent David Gómez-Castro, \\
Departamento de Matem\'{a}ticas, \\
Universidad Aut\'{o}noma de Madrid,\\ Campus de Cantoblanco, 28049 Madrid, Spain.  \\
e-mail address:~\texttt{david.gomezcastro@uam.es}\\
webpage:  \url{https://gomezcastro.xyz/}

\bigskip 

\noindent \L{}ukasz P\l{}ociniczak, \\ 
Department of Applied Mathematics, \\ 
Faculty of Pure and Applied Mathematics,\\ 
Wrocław University of Science and Technology\\
Hugo Steinhaus Center, 
ul. Wybrzeże Wyspiańskiego 27, 
50-370 Wrocław, 
Poland.  \\
e-mail address:~\texttt{lukasz.plociniczak@pwr.edu.pl}\\
webpage:  \url{https://alfa.im.pwr.edu.pl/~plociniczak/}

\bigskip

\noindent Juan Luis V\'azquez, \\ 
Departamento de Matem\'{a}ticas, \\
Universidad Aut\'{o}noma de Madrid,\\ Campus de Cantoblanco, 28049 Madrid, Spain.  \\
e-mail address:~\texttt{juanluis.vazquez@uam.es}\\
webpage:  \url{https://verso.mat.uam.es/\~juanluis.vazquez/}



\end{document}
